% coding: utf-8
% --------------------------------------------------------------------------------------------------
% "Синтез оптимального стохастического управления", 2011 год
% --------------------------------------------------------------------------------------------------
\documentclass[12pt, a4paper, titlepage, openbib, oneside]{report}

\usepackage[utf8]{inputenc}
\usepackage[russian]{babel}
%\usepackage{indentfirst}
\usepackage{amsmath}
\usepackage{amssymb}
\usepackage{theorem}
\usepackage[russian]{varioref}
\usepackage[pdftex]{graphicx}
\usepackage{listings}

\frenchspacing
\fussy
\raggedbottom

\setlength{\leftmargin}{20mm}
\setlength{\rightmargin}{10mm}
\righthyphenmin=2
\emergencystretch=6pt


\newtheorem{teo}{Теорема}
\newtheorem{lemma}{Лемма}
\newtheorem{alg}{Алгоритм}
\newtheorem{df}{Определение}



\DeclareMathOperator{\diag}{\mathrm{diag}}

%\includeonly{titlepage,abstract,toc}

\begin{document}

\renewcommand{\bibname}{Список использованной литературы}
\renewcommand{\chaptername}{Раздел}

% --------------------------------------------------------------------------------------------------
% команды
% --------------------------------------------------------------------------------------------------

% перевод строки
\newcommand{\br}{\vspace{12pt}}

% равно по определению
\newcommand{\eqdef}{\stackrel{\mathrm{def}}{=}}

% пространство R в степени x
\newcommand{\Rv}[1]{\mathbb{R}^{#1}}

% матричное пространство R в степени x на y
\newcommand{\Rs}[2]{\mathbb{R}^{{#1} \times {#2}}}

% математическое ожидание
\newcommand{\E}{\mathbf{E}}

% эпсилон
\newcommand{\eps}{\varepsilon}

% сделать текст полужирным
\newcommand{\strong}[1]{{\fontseries{bx}\selectfont {#1}}}

% частная производная
\newcommand{\PartDiff}[2]{\frac{\partial {#1}}{\partial {#2}}}

% вторая частная производная
\newcommand{\PartDiffSecond}[3]{\frac{\partial^2 {#1}}{\partial {#2}\,\partial {#3}}}

% квадратик qed
\newcommand{\qed}{\vspace{12pt}\hfill \mbox{\raggedright \rule{.1in}{.1in}}}


% --------------------------------------------------------------------------------------------------
% страницы и главы
% --------------------------------------------------------------------------------------------------



% титульная страница
% (готово)
% coding: utf-8
% --------------------------------------------------------------------------------------------------
% "Численные методы аппроксимации гибридных диффузионных процессов", 2010
% : титульный лист
% --------------------------------------------------------------------------------------------------



\begin{titlepage}

    \begin{center}
        Федеральное агентство по образованию\vspace{0.7mm}
        
        Государственное образовательное учреждение\vspace{0.7mm}
        
        высшего профессионального образования\vspace{1.5mm}
        
        <<Нижегородский Государственный университет им.~Н.\,И.\,Лобачевского>>
        
        \vspace{1.2cm}

        {
            \fontsize{18pt}{18pt} \fontseries{bx} \selectfont
            Синтез оптимального регулятора в стохастических системах
        }

        \vspace{2cm}
        
        {\large Магистерская диссертация}
    \end{center}

    \vspace{2cm}

    {
        \fontsize{10pt}{14pt} \selectfont
        \begin{tabular}[n]{p{3.75cm} s{5.22cm} r}
            Выполнил:                 & \multirow{2}*{}                                           &                 \\
            студент гр.~86~М1         &                                                           & Архипов С.\,В.  \\
            \\
            Научный руководитель:     & \multirow{2}*{\includegraphics[scale=0.65]{pakshin_sign}} &                 \\
            профессор, д.\,ф.--м.\,н. &                                                           & Пакшин П.\,В.   \\
            \\
            Заведующий кафедрой:      & \multirow{2}*{}                                           &                 \\
            профессор, д.\,ф.--м.\,н. &                                                           & Федоткин М.\,А. \\
        \end{tabular}
    }

    \vspace{3cm}

    \begin{center}
    {
        \fontsize{10pt}{3mm} \fontshape{sc} \selectfont
            Нижний Новгород \\
            2011 год
    }
    \end{center}

\end{titlepage}



% краткое описание рассматриваемой темы
% coding: utf-8
% --------------------------------------------------------------------------------------------------
% "Численные методы аппроксимации гибридных диффузионных процессов", 2010
% : аннотация
% --------------------------------------------------------------------------------------------------

\begin{abstract}

В работе исследуются гибридные диффузионные системы (гибридные переключаемые диффузии), изучается подход, даются примеры. Дается математическая постановка задачи, приведены основные результаты, связанные с регулярностью, слабой непрерывностью, соотношением с феллеровыми процессами и непрерывной зависимостью от начальных условий.

Приведен один из численных методов решения поставленной задачи, изучен характер сходимости, исследован вопрос существования и единственности решения.

Приведено описание разработанного программного комплекса, реализующего рассчеты динамики гибридных диффузионных систем с независимыми от состояния переключаемыми процессами, дается руководство пользователя, приведены примеры решения поставленных задач с анализом полученных результатов.

\end{abstract}



% содержание. генерируется автоматически
% coding: utf-8
% --------------------------------------------------------------------------------------------------
% "Численные методы аппроксимации гибридных диффузионных процессов", 2010
% : содержание
% --------------------------------------------------------------------------------------------------

\setcounter{page}{1}
\tableofcontents


% введение в тему
% coding: utf-8
% --------------------------------------------------------------------------------------------------
% "Синтез оптимального стохастического управления", 2011 год
% --------------------------------------------------------------------------------------------------



\chapter*{Введение}
\addcontentsline{toc}{chapter}{Введение}



Процессы управления происходят повсюду: и в живой, и в неживой природе: в биологических организмах, обществе, технике. Даже естественный отбор, являющийся основой учения Дарвина, благодаря которому одни особи исчезают, а другие выживают и воспроизводятся, также является своего рода процессом управления, протекающим в природе.

Общим для всех процессов управления, где бы они не протекали, является прием, хранение, преобразование информации и выработка, синтез, управления на ее основе. Осознание этой общности послужило предпосылкой к возникновению в конце сороковых годов XX века научного направления, названного Н.\,Винером \emph{кибернетикой}. Хотя управление человеческими коллективами с одной стороны, и техническими объектами~--- с другой, имеет много общего, коренные различия, которые существуют между этими <<объектами>>, делают необходимыми их раздельное рассмотрение.

В технике \emph{управлением} называют целенаправленное воздействие на какое-либо устройство или объект. Если управление осуществляет человек, то управление называют \emph{ручным}; \emph{автоматическое} управление осуществляется без непосредственного участия человека. Устройство (машина, агрегат, технологический процесс), состоянием которого нужно управлять, называется \emph{объектом управления} или \emph{управляемым объектом}. Целью управления таким объектом является поддержание заданного режима работы или вывод объекта на такой режим работы. Под \emph{заданным режимом} понимают изменение какого-либо параметра, характеризующего состояние объекта управления, по определенному закону.

Объект управления с взаимодействующим с ним управляющим устройством называют \emph{системой управления}. В простейших случаях систему автоматического управления называют \emph{системой автоматического регулирования} (САР), управляющее устройство~--- \emph{регулятором}, а объект управления~--- \emph{объектом регулирования} или \emph{регулируемым объектом}.

\br

Выделяют три основных принципа управления:

\bdescr
	\item[Принцип программного управления]
		Если об объекте точно известно, как зависит выходная переменная от управляющего воздействия, управление можно формировать как известную функцию времени. Такой принцип неприменим при управлении объектом, на который действуют заранее не известные возмущения, оказывающие существенное влияние на управляемую величину. Он также неприменим, если объект является нейтральным или неустойчивым, а система должна функционировать длительное время.
		
	\item[Принцип компенсации]
		При таком принципе управления определяются каким-либо образом действующие на систему управления возмущения и на их основе вырабатывается управляющее воздействие, которое полностью или частично компенсирует влияние возмущений на процесс управления. Также этот принцип называют \emph{способ управления по возмущению} или \emph{принципом Понселе}. Достоинством такого способа управления является принципиальная возмжность полной компенсации возмущающего воздействия.
		
	\item[Принцип обратной связи]
		Этот принцип также называется \emph{управлением по отклонению}. Это такой способ управления, при котором определяется отклонение текущего значения выходной переменной от требуемого значения, и на его основе формируется управляющее воздействие. Системы управления, основанные на этом способе, непременно содержат \emph{обратную связь}~--- канал связи, по которому информация об управляемой переменной поступает на управляющее устройство. Иногда этот принцип называют \emph{принципом Ползунова--Уатта}.
		
		Недостатком такого способа управления является принципиальная невозможность полной компенсации возмущающих воздействий. Это связано с тем, что при этом способе управляющее воздействие начинает вырабатываться и оказывать влияние на ход процесса только после того, как возмущение, начав действовать, вызывает отклонение управляемой величины от требуемого режима.
		
	\item[Принцип комбинированного управления]
		Этот принцип является попыткой совместить достоинства управлений по возмущению и по отклонению. Этот принцип используется в тех случаях, когда на систему действует много различных возмущений, но лишь немногие оказывают наибольшее влияние на работу системы уравления. В подобных случаях влияние превалирующего возмущения можно нейтрализовать, используя принцип компенсации, и нейтрализовать влияние остальных возмущений, используя принцип обратной связи.
\edescr

Задачей работы является изучение теоретической возможности управления объектами, используя принцип комбинированного управления, при учете множества различных случайных возмущений. Это могут быть как шум среды, так, например, механические воздействия друг на друга сразу нескольких объектов, принимающих участие в сложном движении. В результате работы выработаны некоторые алгоритмы, которые могут быть использованы при конструировании требуемого регулятора, отмечены как положительные стороны алгоритмов, так и отрицательные. Каждый алгоритм имеет под собой математическую основу, гарантирующую, что сконструированный таким образом регулятор будет обеспечивать требуемое в определенном смысле управление. Требуемое управление, иначе называемое \emph{оптимальным} говорит о том, что его использование позволит получить требуемый результат наилучшим среди прочих образом.

В процессе исследований был разработан некоторый математический пакет, позволяющий программно сконструировать подобный регулятор и проводить различные экспериментальные проверки с целью выявляения особенностей работы того или иного алгоритма. Пакет распространяется свободно, на принципах открытого сообщества.

\br

Работа устроена следующим образом:

\bdescr
	\item[Постановка задачи]
		Раздел посвящен математической формализации задачи синтеза оптимального регулятора. Дается как математическое описание системы (в случаях случайных воздействий и их отсутствия), так и принцип, по которому определяется оптимальность управления.
	
	\item[Синтез оптимального управления]
		Раздел посвящен методам синтеза оптимального управления в различных случаях без случайных воздействий. Дается математическое обоснование полученных результатов, как в общем случае, так и в частном, наиболее распространенном случае с квадратичным критерием качества. Результаты связываются с принципом Беллмана и принципом максимума.
		
	\item[Методы стохастического управления]
		Раздел посвящен обобщению результатов, полученных в разделе <<\emph{Синтез оптимального управления}>> на случаях стохастических систем, то есть систем, где учитываются случайные воздействия. Дополнительно раскрывается вопрос стохастической устойчивости полученных теоретических результатов.
		
	\item[Пример решения типовой задачи]
		Раздел посвящен решению конкретной типовой задачи, в которой используются результаты из раздела <<\emph{Методы стохастического управления}>>. Дается описание задачи (полет баллистической ракеты в турбулентной атмосфере), конструируются алгоритмы ее решения, дается описание созданного для ее исследования программного математического пакета, анализируются полученные результаты, дается их обоснование.
\edescr


% математическая постановка задачи.
% coding: utf-8
% --------------------------------------------------------------------------------------------------
% "Синтез оптимального стохастического управления", 2011 год
% --------------------------------------------------------------------------------------------------



\chapter{Постановка задачи}
% ==============================================================================================
\newcommand{\funcF}{ \calf{F}         } % функция качества
\newcommand{\optF}{  \optimum{\funcF} } % оптимальная функция F
% ==============================================================================================



Предположим, что имеется некоторый физический объект, чья динамическое состояние описывается с помощью вектора параметров, полностью определяющих его положение в некотором пространстве и достаточных для того, чтобы однозначно определить его положение в следующий момент времени. Динамику его перемещений можно описать с помощью системы дифференциальных уравнений. Не нарушая общности положим, что динамика описывается с помощью системы линейных дифференциальных уравнений

\beq{eq:1/1}
    \eqsystem{
        \dot{x}_1 = a_{11}(t)x_1(t) + a_{12}(t)x_2(t) + \cdots + a_{1n}(t)x_n(t) \text {,} \\
        \dot{x}_2 = a_{21}(t)x_1(t) + a_{22}(t)x_2(t) + \cdots + a_{2n}(t)x_n(t) \text {,} \\
        \vdots                                                                             \\
        \dot{x}_n = a_{n1}(t)x_1(t) + a_{n2}(t)x_2(t) + \cdots + a_{nn}(t)x_n(t) \text {.} \\
    }
\eeq

В матричном виде уравнение~\ref{eq:1/1} можно переписать следующим образом:

\beq{eq:1/2}
    \dot{\m{x}}(t) = \m{A}(t)\m{x}(t) \text{,}
\eeq

где $\m{x}(t) \eqdef (x_1(t), x_2(t), \ldots, x_n(t)) \in \Rv{n}$, а матрица $\m{A}(t) \in \Rs{n}{n}$ и представляется как

\beqn
    \m{A}(t) \eqdef \matr{
        a_{11}(t) & a_{12}(t) & \cdots & a_{1n}(t) \\
        a_{21}(t) & a_{22}(t) & \cdots & a_{2n}(t) \\
        \vdots    & \vdots    & \ddots & \vdots    \\
        a_{n1}(t) & a_{n2}(t) & \cdots & a_{nn}(t)
    } \text{.}
\eeqn

Предположим также, что имеется некоторый способ воздействовать на динамику системы, чтобы было возможным перевести ее из начального состояния $t_0$ в конечное состояние $T$ так, что $\m{x}(T)$ удовлетворяла бы ряду заранее заданных условий. Такие воздействия называются \emph{управлением} и моделируются с помощью вектора $\m{u}(t) \eqdef (u_1(t), u_2(t), \ldots u_m(t)) \in \Rv{m}$.

Вводя управление, систему~\ref{eq:1/2} можно переписать следующим образом:

\beq{eq:1/3}
    \dot{\m{x}}(t) = \m{A}(t)\m{x}(t) + \m{B}(t)\m{u}(t) \text{,}
\eeq

где матрица $\m{B}(t) \in \Rs{n}{m}$ и представляется как

\beqn
    \m{B}(t) \eqdef \matr{
        b_{11}(t) & b_{12}(t) & \cdots & b_{1m}(t) \\
        b_{21}(t) & b_{22}(t) & \cdots & b_{2m}(t) \\
        \vdots    & \vdots    & \ddots & \vdots    \\
        b_{n1}(t) & b_{n2}(t) & \cdots & b_{nm}(t)
    } \text{.}
\eeqn

В работе, не нарушая общности, будут рассматриваться матрицы $\m{B}(t)$ с постоянными значениями, то есть $\m{B}(t) \equiv \m{B}$. Такие предположения сделаны лишь для того, чтобы упростить выкладки.

Используя различное управление, система будет переходить в различные конечные состояния $\m{x}(T)$: используя управление $\m{u}^1(t)$, система будет переведена в конечное состояние $\m{x}^1(t)$, управление $\m{u}^2(t)$~--- в $\m{x}^2(T)$ и так далее. Как правило, имеет значение, в котором окажется система в конечный момент времени. Управление\footnote{Конечно же, оптимальное управление в общем случае не является единственным.}, которое переводит систему в требуемое состояние $\optimum{\m{x}}(t)$ будет называться \emph{оптимальным} и обозначаться соответственно $\optimum{\m{u}}(t)$. Траекторию изменения состояний от момента времени $t_0$ до $T$ также назовем \emph{оптимальной}.

\br

Зачастую имеет значение не только факт, что в определенный момент времени система перешла в нужное состояние, но и то, как был совершен этот переход. Иными словами, требуется учитывать предысторию перехода, всю динамику системы с использованием управлений. Чтобы учитывать такое <<качество>> перехода, вводятся различные \emph{критерии качества}~--- функции, выражающие качество перехода через некоторые количественные значения. Так как для вещественных чисел определены операции сравнения, то появляется возможность таким образом сравнивать траектории и выделять оптимальное уже в этом случае управление. Иногда даже не имеет значения, в каком конечном состоянии окажется система, важно, чтобы переход осуществлялся так, чтобы он удовлетворял некоторому критерию качества сам по себе.

Естественным образом возникает задача минимизации (или максимизации) такого критерия качества $\funcF\bigl(\m{x}(t), \m{u}(t)\bigr) \colon \Rv{n} \times \Rv{m} \to \Rv{1}$. Для определенности рассмотрим задачу минимизации критерия. Минимальное значение $\min\funcF\bigl(\m{x}(t), \m{u}(t)\bigr)$ обозначим как $\optF\bigl(\m{x}(t), \m{u}(t)\bigr)$.

При этих предположениях, задачу можно сформулировать таким образом: требется найти управление $\optimum{\m{u}}(t)$ для системы~\ref{eq:1/2} и заданного критерия качества $\funcF$ при условии

\beq{eq:1/4}
    \funcF\bigl(\m{x}(t), \m{u}(t)\bigr) \to \min \text{.}
\eeq

Найденное таким образом управление также будет называться оптимальным.

\br

Однако система~\ref{eq:1/3} помимо воздействия управления может подвергаться и воздействиям различных внешних факторов, которые нельзя проконтроллировать. Воздействие этих факторов непредсказуемо, они оказывают влияние на поведение объекта на протяжение всего интервала наблюдений $[t_0, T]$. Такие случайные воздействия возникают, например, вследствие разнообразных шумов среды, столкновений с другими объектами или скачков энергии. Факторов может быть огромное количество, но при моделировании их можно объединять в группы. Эти воздействия описываются случайными величинами, в частности векторными. В работе будет исследоваться влияние векторной случайной величины $\m{\xi}(t) = \bigl( \xi_1(t), \xi_2(t), \ldots, \xi_k(t) \bigr) \in \Rv{k}$.

Таким образом, задачу~\ref{eq:1/3}--\ref{eq:1/4} можно переписать следующим образом:

\beq{eq:1/5}
    \dot{\m{x}}(t) = \m{A}(t)\m{x}(t) + \m{B}(t)\m{u}(t) + \m{\Sigma}(t)\m{\xi}(t) \text{,}
\eeq

где $\m{\Sigma}(t) \in \Rs{n}{k}$ и записывается как

\beqn
    \m{\Sigma}(t) \eqdef \matr{
        \sigma_{11}(t) & \sigma_{12}(t) & \cdots & \sigma_{1k}(t) \\
        \sigma_{21}(t) & \sigma_{22}(t) & \cdots & \sigma_{2k}(t) \\
        \vdots    & \vdots    & \ddots & \vdots    \\
        \sigma_{n1}(t) & \sigma_{n2}(t) & \cdots & \sigma_{nk}(t)
    } \text{.}
\eeqn

Впрочем, не нарушая общности, как и в случае для матрицы $\m{B}(t)$ будем полагать, что матрица $\m{\Sigma}(t)$ является постоянной, то есть $\m{\Sigma}(t) \equiv \m{\Sigma}$. Критерий качества, очевидно, остается прежним, поскольку факторы, которые нельзя проконтроллировать, влияют на саму траекторию системы.

Нахождение оптимальных управлений $\optimum{\m{u}}(t)$ для системы~\ref{eq:1/5} и задачи минимизации~\ref{eq:1/4} и является задачей этой работы.

% синтез оптимального управления
\chapter{Синтез оптимального управления}
% ==============================================================================================
\newcommand{\optF}{  \optimum{\calf{F}} } % оптимальная функция F
\newcommand{\optU}{  \optimum{u}        } % оптимальное управление u
\newcommand{\optX}{  \optimum{x}        } % оптимальная траектория x
\newcommand{\funcL}{ \sff{L}            } % функция Лагранжа
% ==============================================================================================




% **********************************************************************************************
\section{Принцип Беллмана}
% **********************************************************************************************



Рассмотрим следующую задачу Лагранжа. Для системы $\dot{x} = f(x, u, t)$, начинающей движение из начального состояния $x(t_1) = x_1$, найдем управление $u(t)$, ограниченное некоторым допустимым классом функций $\Omega(t)$ и определенное на интервале времени $[t_1, t_2]$, которое минимизирует показатель качества

\beq{eq:2/1/1}
    \calf{F} = \int\limits_{t_1}^{t_2} \funcL(x, u, t)\,dt \mbox{,}
\eeq

где функция $\funcL(x, u, t)$ предполагается непрерывной по $t$. Конечное состояние $x(t_2) = x_2$ не задано. Все функции, принадлежащие $\Omega(t)$ ограничены по амплитуде в любой момент времени $t$: $|u_j(t)| \leqslant V_j$. Множество точек $|u_j(t)| \leqslant V_j$ ($j = 1, 2, \ldots, n$) будем обозначать как $\Upsilon$.

Показатель качества $\calf{F}$ для фиксированного значения $t_2$ зависит от переменных $u(t)$, $x_1$ и $t_1$. Однако оптимальное значение функционала $\calf{F}$ зависит лишь от начального состояния $x_1$ и момента времени $t_1$. Обозначим оптимальное значение этого функционала как $\optF(x_1, t_1)$. Если поставлена задача оптимального управления, то для каждой точки $x_1$ ставится в соответствие некоторое значение критерия оптимальности. Функция $\optF(x_1, t_1)$ определяет гиперповерхность в $n+1$-мерном пространстве. Эта гиперповерхность, конечно, в явном виде не задана, но некоторые ее свойства позволяют вывести ряд необходимых условий, которым должна удовлетворять функция оптимального управления $\optU(t)$.

Обозначим через $x \bigl( u(t), t \bigr)$ траекторию, которая получается в результате воздействия управления $u(t)$ на систему $\dot{x} = f(x, u, t)$ с начальным состоянием $x_1$ в момент времени $t = t_1$. При этом $\optF$ определяется выражением

\beq{eq:2/1/2}
    \optF(x_1, t_1) = \underset{\substack{u(t) \in \Omega(t) \\ t_1 \leqslant t \leqslant t_2}}{\min} \biggl\{ \int\limits_{t_1}^{t_2} \funcL\bigl(x(u, t), u, t\bigr)\,dt \biggr\} \mbox{.}
\eeq

Для некоторого момента времени $t'$ из интервала между $t_1$ и $t_2$ это выражение можно написать в таком виде:

\beq{eq:2/1/3}
    \optF(x_1, t_1) = \underset{\substack{u(t) \in \Omega(t) \\ t_1 \leqslant t \leqslant t_2}}{\min} \biggl\{ \int\limits_{t_1}^{t'} \funcL\bigl(x(u, t), u, t\bigr)\,dt + \int\limits_{t'}^{t_2} \funcL\bigl(x(u, t), u, t\bigr)\,dt \biggr\} \mbox{.}
\eeq

Уравнение~\ref{eq:2/1/3} позволяет применить для его решения принцип оптимальности. Для случая непрерывных систем принцип оптимальности можно сформулировать следующим образом.

\begin{statement}\lbl{statement:1}
	Оптимальное управление $\optU(t)$ на интервале времени $[t_1, t_2]$ имеет следующее свойство: для любого $t'$, заключенного в интервале $t_1 < t' < t_2$ независимо от значений, которые управление $\optU(t)$ принимало на интервале времени $[t_1, t']$, и, следовательно, независимо от значения $\optX(t')$ оно должно оставаться оптимальным управлением относительно состояния $\optX(t')$ на интервале времени $(t', t_2]$.
\end{statement}

Применяя этот принцип, уравнение~\ref{eq:2/1/3} можно преобразовать к следующему виду

\beq{eq:2/1/4}
    \optF(x_1, t_1) = \underset{\substack{u(t) \in \Omega(t) \\ t_1 \leqslant t \leqslant t'}}{\min} \biggl\{ \int\limits_{t_1}^{t'} \funcL\bigl(x(u, t), u, t\bigr)\,dt + \optF\bigl(x(t'), t'\bigr) \biggr\} \mbox{,}
\eeq

где $x(t')$~--- конечное состояние, которое является результатом действия управления $u(t)$ на интервале времени $[t_1, t']$\footnote{ Если предположить, что функция $\calf{F}$ имеет непрерывные частные производные по $x$ и $t$, можно легко прийти к уравнению Беллмана (см TODO). Однако, как будет показано ниже, это предположение не выполняется для большого класса задач оптимального управления. В приводимом здесь доказательстве, основанном на работе TODO, это предположение не используется }.

При оптимальном управлении $u(t) = \optU(t)$ на интервале $[t_1, t']$ имеем

\beq{eq:2/1/5}
    \optF(x_1, t_1) = \int\limits_{t_1}^{t'} \funcL\bigl(x(\optU, t), \optU(t), t\bigr)\,dt + \optF\bigl(x(t'), t'\bigr) \mbox{.}
\eeq

Перенося члены и деля на $t' - t_1$ получим

\beq{eq:2/1/6}
    -\frac{\optF\bigl( x(t'), t' \bigr) - \optF( x_1, t_1)}{t' - t_1} = \frac{1}{t' - t_1} \int\limits_{t_1}^{t'} \funcL\bigl(x(\optU, t), \optU(t), t\bigr)\,dt \mbox{.}
\eeq

При $t' \to t_1$ уравнение~\ref{eq:2/1/6} принимает вид

\beq{eq:2/1/7}
    \lim_{t' \to t_1} \biggr[ -\frac{\optF\bigl( x(t'), t' \bigr) - \optF( x_1, t_1)}{t' - t_1} \biggr] = \funcL\bigl(x_1, \optU(t_1), t_1\bigr) \mbox{.}
\eeq

Заметим, что для получения правой части выражения~\ref{eq:2/1/7} как предела правой части уравнения~\ref{eq:2/1/6} используется теорема о среднем значении (TODO). Если предел левой части уравнения~\ref{eq:2/1/7} существует, то можно определить величину

\beq{eq:2/1/8}
    \biggl[ \genericdiff{\optF}{t} \biggr]_{\optU, t_1} \eqdef \lim_{t' \to t_1} \biggr[ -\frac{\optF\bigl( x(t'), t' \bigr) - \optF( x_1, t_1)}{t' - t_1} \biggr] \mbox{.}
\eeq

Величина $[ \genericlinediff{\optF}{t} ]_{\optU, t_1}$ есть производная функции $\optF$ по времени, вычисленная в момент времени $t_1$. Анализируя правую часть выражения~\ref{eq:2/1/8}, можно видеть, что производную следует вычислять вдоль траектории, обусловленной управлением $\optU$ и начинающейся в $x_1$. Иными словами, пусть для некоторой задачи функция $\optF(x, t)$ определена для любого начального состояния $x$ и начального момента времени $t$. Пусть $x\bigl( u(t), t \bigr)$ определяет траекторию, обусловленную управлением $u(t)$; тогда вдоль любой из этих траекторий величина $\optF(x, t)$ будет изменяться во времени со скоростью, определяемой величиной $[ \genericlinediff{\optF}{t} ]_{u, t}$. В частности, вдоль оптимальной траектории она будет изменяться со скоростью $[ \genericlinediff{\optF}{t} ]_{\optU, t}$.

Таким образом

\beq{eq:2/1/9}
    \biggl[ \genericdiff{\optF}{t} \biggr]_{\optU, t_1} + \funcL\bigl(x_1, \optU(t_1), t_1\bigr) = 0 \mbox{.}
\eeq

Заметим, что при $u(t) \neq \optU(t)$ для интеграла~\ref{eq:2/1/5} в соответствии с определением должно выполняться неравенство

\beq{eq:2/1/10}
    \optF(x_1, t_1) \leqslant \int\limits_{t_1}^{t'} \funcL\bigl(x(\optU, t), \optU(t), t\bigr)\,dt + \optF\bigl(x(t'), t'\bigr) \mbox{.}
\eeq

Откуда следует, что

\beq{eq:2/1/11}
    \biggl[ \genericdiff{\optF}{t} \biggr]_{u, t_1} + \funcL\bigl(x_1, u(t_1), t_1\bigr) \geqslant 0 \mbox{.}
\eeq

Используя выражения~\vref{eq:2/1/9} и~\ref{eq:2/1/11}, получим

\beq{eq:2/1/12}
    \biggl[ \genericdiff{\optF}{t} \biggr]_{u, t_1} + \funcL\bigl(x_1, u(t_1), t_1\bigr) \geqslant \biggl[ \genericdiff{\optF}{t} \biggr]_{\optU, t_1} + \funcL\bigl(x_1, \optU(t_1), t_1\bigr) \mbox{,}
\eeq

или

\beq{eq:2/1/13}
    \underset{u(t_1) \in \Upsilon}{\min} \biggl\{ \biggl[ \genericdiff{\optF}{t} \biggr]_{u, t_1} + \funcL\bigl(x_1, u(t_1), t_1\bigr) \biggr\} = 0 \mbox{.}
\eeq

Уравнение~\ref{eq:2/1/13} справедливо в любой момент времени из интервала $[t_1, t_2]$, так что для любого момента времени $t$ и любого состояния $x$ на траектории $\optX(t)$, принимая его за начальное, можно написать

\beq{eq:2/1/14}
    \underset{u(t) \in \Upsilon}{\min} \biggl\{ \biggl[ \genericdiff{\optF}{t} \biggr]_{u, t} + \funcL\bigl(x(u, t), u, t\bigr) \biggr\} = 0 \mbox{.}
\eeq

Уравнение~\ref{eq:2/1/14} и есть функциональное уравнение Беллмана в общей форме, выражающее необходимое условие оптимальности.

Для любого состояния $x$ и момента времени $t$, когда функция $\genericlinediff{\optF}{t}$, $\genericlinediff{\optF}{x}$ и $f(x, u, t)$ непрерывны по $x$ и $t$, полную производную $[\genericlinediff{\optF}{t}]_{u, t}$ можно записать так:

\beq{eq:2/1/15}
    \biggl[ \genericdiff{\optF}{t} \biggr]_{u, t} = \biggl(\partdiff{\optF}{x}\biggr)^T f(x, u, t) + \partdiff{\optF}{t} \mbox{.}
\eeq

Так как частная производная $\partlinediff{\optF}{t}$ не зависит от $u$, уравнение \ref{eq:2/1/14} можно представить в виде

\beq{eq:2/1/16}
    - \partdiff{\optF}{t} = \underset{u(t) \in \Upsilon}{\min} \Biggl\{ \funcL\bigl(x(u, t), u, t\bigr) + \biggl(\partdiff{\optF}{x}\biggr)^T f(x, u, t) \Biggr\} \mbox{.}
\eeq

В тех случаях, когда это уравнение применимо, оно обеспечивает необходимое условие оптимальности. Отметим, что оно является необычной формой дифференциального уравнения в частных производных, которое включает операцию минимизации. Это, в общем случае нелинейное, дифференциальное уравнение в частных производных первого порядка относительно одной переменной $\optF$. Это уравнение определяет одноточечную краевую задачу с граничным условием вида

\beq{eq:2/1/17}
    \lim_{t_1 \to t_2} \optF(x_1, t_1) = \lim_{t_1 \to t_2} \int\limits_{t_1}^{t_2} \funcL(\optX, \optU, t)\,dt = 0 \mbox{.}
\eeq

Очевидно, это условие соответствует двум.

\benum
    \item
        Оно указывает, как следует вести поиск оптимальной $u(t)$: в любой момент времени $t$ поиск должен быть таким, чтобы минимизировать величину, заключенную в квадратные скобки в уравнении~\vref{eq:2/1/16}.

    \item
        Если функция оптимального управления $\optU(t)$ найдена, то~\ref{eq:2/1/16} сводится к уравнению без операции минимизации
        \beq{eq:2/1/18}
            - \partdiff{\optF}{t} = \funcL\bigl(x(\optU(t), t), \optU, t\bigr) + \biggl(\partdiff{\optF}{x}\biggr)^T f(x(\optU(t), t), \optU, t) \mbox{,}
        \eeq 
        которому удовлетворяет функция $\optF(x, t)$ для всех значений $t$ в интервале $[t_1, t_2]$.
\eenum

При отсутствии ограничений на величину $u$ и при условии, что функции $\funcL$ и $f$ имеют частные производные по $u$, оптимальное управление $\optU(t)$ можно найти путем дифференцирования выражения, заключенного в квадратные скобки в уравнении~\ref{eq:2/1/16} и приравнивания полученного результата к нулю. Это дает условие

\beq{eq:2/1/19}
    \biggl[ \partdiff{\funcL}{u} + \sum\limits_{i=1}^n \partdiff{\optF}{x_i} \centerdot \partdiff{f_i}{u} \biggr] = 0 \mbox{ для всех } t \in [t_1, t_2] \mbox{.}
\eeq

Если показатель качества соответствует задаче Майера, а именно $\calf{F} = P\bigl( x(t_2), t_2 \bigr)$, то, придерживаясь тех же рассуждений, получим соответствующее уравнение Беллмана

\beq{eq:2/1/20}
    \underset{u(t) \in \Upsilon}{\min} \Biggl\{ \biggl[ \genericdiff{\optF}{t} \biggr]_{u, t} \Biggr\} = 0
\eeq

с граничным условием

\beq{eq:2/1/21}
    \lim_{t_1 \to t_2} \optF\bigl(x(t_1), t_1\bigr) = P\bigl(x(t_2), t_2\bigr) \mbox{.}
\eeq

В тех же случаях, когда функции $\partlinediff{\optF}{t}$, $\partlinediff{\optF}{x}$ и $f(x, u, t)$ непрерывны по $x$ и $t$, имеем

\beq{eq:2/1/22}
    - \partdiff{\optF}{t} = \underset{u(t) \in \Upsilon}{\min} \Biggl\{ \biggl(\partdiff{\optF}{x}\biggr)^T f(x, u, t) \Biggr\}
\eeq

с тем же граничным условием~\ref{eq:2/1/21}.

Отметим, что тип граничного условия для конечного состояния не играет роли при выводе функционального уравнения. Таким образом, уравнение Беллмана~\ref{eq:2/1/14} или~\vref{eq:2/1/16} справедливо для задачи, даже если конечное состояние $x_2$ задано, а конечный момент времени $t_2$ нет. Условие~\vref{eq:2/1/17} по-прежнему справедливо, и нет необходимости вносить какие-либо изменения при выводе уравнения Беллмана.



% **********************************************************************************************
\section{Синтез оптимального управления в системах с квадратичным критерием качества}
% **********************************************************************************************



% **********************************************************************************************
\section{Связь с принципом максимума}
% **********************************************************************************************



% **********************************************************************************************
\section{Достаточное условие оптимальности}
% **********************************************************************************************


% методы стохастического управления
% coding: utf-8
% --------------------------------------------------------------------------------------------------
% "Синтез оптимального стохастического управления", 2011 год
% --------------------------------------------------------------------------------------------------



\chapter{Методы стохастического управления}
% ==============================================================================================
\newcommand{\mSig}{     \sff{\Sigma}                     } % матрица ковариаций

\newcommand{\setTime}{  [0, T]                           } % временной промежуток, который достало вводить
\newcommand{\setX}{     \calf{X}                         } % множество X
\newcommand{\setT}{     \calf{T}                         } % множество T
\newcommand{\setF}{     \calf{F}                         } % множество F
\newcommand{\setM}{     \calf{M}_2\bigl( \setTime \bigr) } % множество M2
\newcommand{\setK}{     \calf{K}_2\bigl( \setTime \bigr) } % множество K2
\newcommand{\setU}{     \calf{U}                         } % множество U
\newcommand{\setC}{     \calf{C}                         } % множество U

\newcommand{\argEmpty}{ \{\emptyarg\}                    } % пустой аргумент

\renewcommand{\funcF}{  \calf{F}                         } % функция качества
\renewcommand{\funcH}{  \calf{H}                         } % функция качества H
\renewcommand{\funcV}{  \sff{V}                          } % функция Ляпунова

\newcommand{\idxB}[1]{  ^{(#1)}                          } % индекс в скобках
% ==============================================================================================



% **********************************************************************************************
\section{Некоторые сведения из теории случайных процессов}
% **********************************************************************************************



Перед рассмотрением методов стохастического управления, дадим некоторые сведения из теории вероятностей, в частности из теории случайных процессов.

\begin{df}\lbl{df:1}
    Семейство случайных величин, параметризованное параметров $t$, интерпретируемым как время, будем называть случайным процессом и обозначать как
    
    \beqn
    	\xi_t \eqdef \xi(t, \omega) \in \setX \subset \Rv{1} \text{,}
    \eeqn

    где $t \in \setT \subset \Rv{1}$~--- временной параметр; $\setT$~--- множество определения, а $\setX$~--- множество значений случайного процесса $\xi_t$.
\end{df}

Область определения $\setT$ случайного процесса может быть конечным, полубесконечным или бесконечным интервалом числовой прямой. В этих случаях $\xi(t, \omega)$ называется \emph{случайным процессом с непрерывным временем}. Если $\setT$ имеет конечное или счетное множество элементов $t_k \in \setT: t_k < t_{k+1}$ ($k = 0, 1, 2, \ldots$), то $\xi_t$ называется \emph{случайным процессом с дискретным временем} или \emph{случайной последовательностью}.

Случайный процесс считается полностью заданным, если заданы его конечномерные распределения~--- набор функций распределения.

Случайный процесс называется \emph{гауссовским}, если все его совместные плотности распределения являются гауссовскими:

\beqn
	\prob_\xi(x1, x2, \ldots, x_k, t_1, t_2, \ldots, t_k) = \frac{1}{\sqrt{(2\pi)^n \det(\mSig)}}\exp\biggl( -\frac{(x-m)^T\mSig^{-1}(x-m)}{2} \biggr) \text{,}
\eeqn

где $x = (x_1, x_2, \ldots, x_k)^T$; $m = (m_1, m_2, \ldots, m_k)^T$; $m_i = \E x_i$; $\mSig = \mSig^T > 0$~--- матрица ковариаций.

Случайный процесс $\xi_t$, заданный на $\setT$, называется \emph{стационарным в узком смысле}, если для любых $n \geqslant 1$ и любых моментов $t_1, t_2, \ldots, t_n, t$, таких, что $t_i + t \in \setT$ ($i = 1, 2, \ldots, n$), совместные функции распределения совокупности случайных величин $\xi_{t_1+t}, \xi_{t_2+t}, \ldots, \xi_{t_n+1}$ не зависят от $t$.

\begin{df}\lbl{df:2}
	Случайный процесс $\xi_t$, определенный при всех $t \in \setTime$, называется непрерывным в момент времени $t^{*} \in \setTime$ в вероятностном смысле, если для любой последовательности $\{t_k\}_{k=1}^\infty$, принадлежащей $[0, T]$ и сходящейся к $t^{*}$ при $k \to \infty$, справедливо предельное соотношение
	
	\beqn
		\xi_{t_k} \to \xi_{t^{*}} \text{ при } t_k \to t^{*}  \text{,}
	\eeqn
	
	понимаемое в соответствующем вероятностном смысле.
\end{df}

\begin{df}\lbl{df:3}
	Случайный процесс $\xi_t$ называется непрерывным на промежутке $[t_1, t_2]$ в данном вероятностном смысле, если он непрерывен в этом смысле при $\forall t \in [t_1, t_2]$. Непрерывность в среднем степени 2 случайного процесса $\xi_t$ будем называть среднеквадратической непрерывностью.
\end{df}

\begin{df}\lbl{df:4}
	Случайный процесс $\xi_t$, определеный для всех $t \in \setTime$, называется дифференцируемым в момент $t \in \setTime$ в одном из вероятностных смыслов, если существует левая часть следующего предельного соотношения:
	
	\beqn
		\lim\limits_{\delta \to 0} \frac{\xi_{t+\delta} - \xi_t}{\delta} = \genericdiff{\xi_t}{t} \text{,}
	\eeqn
	
	понимаемого в этом вероятностном смысле. Процесс $\xi_t$ будем называть дифференцируемым на промежутке $[t_1, t_2] \subset \setTime$ в данном вероятностном смысле, если он дифференцируем в этом смысле при всех $t \in [t_1, t_2]$.
\end{df}

\begin{df}\lbl{df:5}
	Случайный процесс $\xi_t$, определеный для всех $t \in \setTime$, называется процессом с независимыми приращениями, если для любых $t_0, t_1, \ldots, t_k \in \setTime$ ($0 \leqslant t_0 < t_1 < \ldots < t_k \leqslant T$), случайные величины $\xi_{t_0}, \xi_{t_1}-\xi_{t_0}, \ldots, \xi_{t_k}-\xi_{t_{k-1}}$ стохастически независимы между собой.
\end{df}

Особую роль в дальнейшем изложении будут играть так называемые винеровские случайные процессы, относящиеся к классу гауссовских процессов с независимыми приращениями\footnote{ Подробное изложение теории случайных процессов с независимыми приращениями можно найти в монографии А.\,В.\,Скорохода\cite{SKOROHOD} }.

\begin{df}\lbl{df:6}
	Случайный процесс $\xi_t \in \Rv{1}$, определеный для всех $t \in \setTime$, называется винеровским, если он является процессом с независимыми приращениями и удовлетворяет следующим условиям
	
	\benum
		\item
			$\E \xi_t = 0$,
		
		\item
			$\E (\xi_{t_2} - \xi_{t_1})^2 - \sigma_\xi^2 |t_2 - t_1| \text{, где } \sigma_\xi^2 > 0$,
		
		\item
			$\xi_{t_1}$ и $\xi_{t_2} - \xi_{t_1}$ ($t_2 > t_1$) имеют нормальное распредление при всех $t_1, t_2 \in \setTime$.
	\eenum
\end{df}

Если $\sigma_\xi^2$ и $\xi_0 = 0$, то винеровский процесс называется \emph{стандартным}. Поскольку приращения $\xi_{t_2} - \xi_{t_1}$ ($t_2 > t_1$) винеровского процесса имеют нормальное распределение, то их функция распределения определяется соотношением

\beqn
	\prob(\xi_{t_2} - \xi_{t_1} < x) = \frac{1}{\sigma_\xi \sqrt{2\pi (t_2 - t_1)}} \int\limits_\infty^x e^{-\frac{u^2}{2\sigma_xi^2(t_2-t_1)}}\,du \text{.}
\eeqn

Известно, что винеровский случайный процесс является непрерывным с вероятностью 1 случайным процессом.

\br

Определим теперь понятия интегралов и процессов Ито.

Пусть задано вероятностное пространство $(\Omega, \setF, \prob)$ и винеровский случайный процесс $f_t \in \Rv{1}$. Рассмотрим совокупность $\sigma$-алгебр $\{ \setF_t, t \in \setTime \}$, определенную на вероятностном пространстве $(\Omega, \setF, \prob)$ и связанную с винеровским процессом $f_t$ так, что

\benum
	\item
		$\setF_s \subset \setF_t \subset \setF$ при $s<t$;
	
	\item
		$f_t$ измерим относительно $\setF_t$ при каждом $t \in \setTime$;
		
	\item
		Процесс $f_{t+\Delta} - f_\Delta$ при всех $\Delta \geqslant 0$ и $t > 0$ не зависит от событий $\sigma$-алгебры $\setF_\Delta$.
\eenum

В дальнейшем рассматриваются только такие винеровские процессы $f_t$, которые удовлетворяют условиям 2 и 3.

Введем в рассмотрение класс $\setM$ функций $\xi: \setTime \times \Omega \to \Rv{1}$, которые удовлетворяют следующим условиям:

\benum
	\item
		Функции $\xi(t, \omega)$ измеримы по совокупности переменных $(t, \omega)$;
	
	\item
		Функция $\xi(t, \omega) \in \setF_t$ измерима при всех $t \in \setTime$, и ее значения $\xi(\tau, \omega)$ не зависят от приращений $f_{t+\Delta} - f_\Delta$ винеровского процесса при $\Delta \geqslant \tau$, $t > 0$;
	
	\item
		$\int_0^T \E \Bigl\{ \bigl( \xi(t, \omega) \bigr)^2 \Bigr\}\,dt < \infty$;
	
	\item
		$\E \Bigl\{ \bigl( \xi(t, \omega) \bigr)^2 \Bigr\} < \infty$ для всех $t \in \setTime$.
\eenum

Введем на классе $\setM$ норму вида

\beqn
	\| \xi \|_{2,T} = \sqrt{ \int_0^T \E \Bigl\{ \bigl( \xi(t, \omega) \bigr)^2 \Bigr\}\,dt } \text{.}
\eeqn

Считаем две функции $\xi$ и $\eta$ из класса $\setM$ эквивалентными, если $\| \xi - \eta \|_{2,T} = 0$.

Для произвольного разбиения $\{ \tau_j \}_{j=0}^N$ промежутка $\setTime$, такого что $0 = \tau_0 < \tau_1 < \ldots < \tau_N = T$, и произвольной последовательности среднеквадратически интегрируемых, $\setF_{\tau_j}$-измеримых и независящим от приращений $f_{s+t} - f_t$ при $t \geqslant \tau_j$, $s > 0$ случайных величин $\xi_j(\omega)$ ($j = 0, 1, \ldots, N-1$), определим \emph{ступенчатую функцию} $\xi \in \setM$ следующим образом: $\xi(t, \omega) = \xi_j(\omega)$ с вероятностью 1 при $\tau_j \leqslant t < \tau_{j+1}$ ($j = 0, 1, \ldots, N-1$).

При этом интеграл в требовании 3 понимается как

\beqn
	\int\limits_0^T \E \Bigl\{ \bigl( \xi(t, \omega) \bigr)^2 \Bigr\}\,dt = \sum\limits_{j=0}^{N-1} \E \bigl\{ (\xi_j)^2 \bigr\}(\tau_{j+1} - \tau{j}) \text{.}
\eeqn

Обозначим через $\setK$ множество всех ступенчатых функций в $\setM$. Для произвольной ступенчатой функции $\xi$ из класса $\setK$ определим \emph{стохастический интеграл Ито} следующим образом:

\beqn
	\int\limits_0^T \xi_\tau\,df_\tau \eqdef \sum\limits_{j=0}^{N-1} \xi_j(\omega)\bigl( f(\tau_{j+1}, \omega) - f(\tau_j, \omega) \bigr) \text{.}
\eeqn

\begin{df}\lbl{df:7}
	Стохастическим интегралом Ито от функции $\xi \in \setM$ назовем среднеквадратический предел последовательности ступенчатых функций:
	
	\beqn
		\int\limits_0^T \xi_\tau\,df_\tau = \lim\limits_{N \to \infty} \sum\limits_{j=0}^{N-1} \xi\idxB{N}\bigl( f(\tau\idxB{N}_{j+1}, \omega) - f(\tau\idxB{N}_j, \omega) \bigr) \text{,} 
	\eeqn
	
	где $\xi\idxB{N}(t, \omega)$~--- произвольная ступенчатая функция из $\setK$, сходящаяся по норме $\| \emptyarg \|_{2,T}$ к функции $\xi(t, \omega)$.
\end{df}

Известно, что для $\xi \in \setM$ такой интеграл существует\cite{KLOEDEN}.

\br

Винеровские случайные процессы порождают широкий класс случайных процессов, называемых процессами Ито.

Пусть задано фиксированное вероятностное пространство $(\Omega, \setF, \prob)$ и $\setF_t$-измеримый при всех $t \in \setTime$ векторный винеровский процесс $\m{f}_t \in \Rv{m}$ с независимыми компонентами $\m{f}_t\idxB{i}$ ($i = 1, 2, \ldots, m$).

Рассмотрим случайные процессы $a_t \in \Rv{1}$ и $\m{b}_t\idxB{i} \in \Rv{1}$ ($i = 1, 2, \ldots, m$) из класса $\setM$. Определим случайный процесс $\eta_t$ при всех $t \in \setTime$ в виде:

\beq{eq:3/1/1}
	\eta_t = \eta_0 + \int\limits_0^t a_\tau\,d\tau + \sum\limits_{i=1}^m \int\limits_0^t \m{b}_\tau\idxB{i}\,d\m{f}_\tau\idxB{i} \text{,}
\eeq

где $\eta_0$~--- случайная величина, стохастически независимая от приращений $\m{f}_\tau-\m{f}_0$ винеровского процесса $\m{f}_t$ при всех $\tau \in \setTime$.

\begin{df}\lbl{df:8}
	Процесс $\eta_t$, определенный на промежутке $\setTime$ в виде~\ref{eq:3/1/1}, называется процессом Ито, если правая часть~\ref{eq:3/1/1} существует при всех $t \in \setTime$. При этом процесс $a_t$ называется коэффициентом сноса, а процессы $\m{b}_t\idxB{i}$ ($i = 1, 2, \ldots, m$)~--- коэффициентами диффузии процесса $\eta_t$.
\end{df}

Из свойства аддитивности стохастических интегралов вытекает следующее представление для процесса Ито:

\beq{eq:3/1/2}
	\eta_s = \eta_t + \int\limits_t^s a_\tau\,d\tau + \sum\limits_{i=1}^m \int\limits_0^t \m{b}_\tau\idxB{i}\,d\m{f}_\tau\idxB{i} \text{ с вероятностью 1,}
\eeq

справедливое для всех $s, t \in \setTime$ и $s \geqslant t$.

Рассмотрим многомерный случай. Пусть существуют случайные процессы $\m{a}_t \in \Rv{n}$ и $\m{B}_t \in \Rs{n}{m}$, каждая компонента которых должна принадлежать классу $\setM$. Определим векторный случайный процесс $\m{\eta}_t \in \Rv{n}$ при всех $t \in \setTime$ в виде:

\beq{eq:3/1/3}
	\m{\eta}_t = \m{\eta}_0 + \int\limits_0^t \m{a}_\tau\,d\tau + \int\limits_0^t \m{B}_\tau\,d\m{f}_\tau \text{,}
\eeq

где случайная векторная величина $\m{\eta}_0 \in \Rv{n}$ является стохастически независимой от приращений $\m{f}_\tau - \m{f}_0$ винеровского процесса $\m{f}_t$ при всех $\tau \in \setTime$.

\begin{df}\lbl{df:9}
	Процесс $\m{\eta}_t \in \Rv{n}$ вида~\ref{eq:3/1/2}, определенный на промежутке $\setTime$ называется векторным процессом Ито, если правая часть~\ref{eq:3/1/2} существует при всех $t \in \setTime$. При этом процессы $\m{a}_t\idxB{i}$ называются коэффициентами снова, а процессы $\m{B}_t\idxB{i}$ называются коэффициентами диффузии процесса $\m{\eta}_t$ ($i = 1, 2, \ldots, n$, $j = 1, 2, \ldots, m$).
\end{df}

Соотношение~\ref{eq:3/1/2} обобщается и на случай векторного процесса Ито:

\beq{eq:3/1/4}
	\eta_s = \eta_t + \int\limits_t^s \m{a}_\tau\,d\tau + \int\limits_0^t \m{B}_\tau\,d\m{f}_\tau \text{ с вероятностью 1,}
\eeq

справедливое для всех $s, t \in \setTime$ и $s \geqslant t$.

\br

Процессы Ито тесно связаны со стохастическими дифференциальными уравнениями Ито. Действительно, пусть задано дифференциальное уравнение

\beq{eq:3/1/5}
	dx_t = a(x_t, t)dt + \sigma(x_t, t)df_t \text{; } x_0 = x(0, \omega) \text{,}
\eeq

где $x_t \in \Rv{1}$~--- случайный процесс, являющийся решением уравнения~\ref{eq:3/1/5}; $a(x, t)$ и $\sigma(x, t) \in Rv{1}$~--- измеримые при всех $(x, t) \in \Rv{1} \times \setTime$ функции. Тогда это уравнение может быть переписано в интегральном виде

\beq{eq:3/1/6}
	x_t = x_0 + \int\limits_0^t a(x_\tau, \tau)\,d\tau + \int\limits_0^t \sigma(x_\tau, \tau)\,df_\tau \text{; } x_0 = x(0, \omega) \text{.}
\eeq

Второй интеграл в правой части~\ref{eq:3/1/6} может пониматься в смысле Ито, Стратоновича или как $\theta$-интеграл\cite{KUZNETCOV}. Пусть интеграл $\int_0^t \sigma(x_\tau, \tau)\,df_\tau$ является стохастическим интегралом Ито, тогда уравнение~\ref{eq:3/1/5} и соответствующее ему уравнение~\ref{eq:3/1/6} называются стохастическими дифференциальными уравнениями Ито.



% **********************************************************************************************
\section{Оптимальное стохастическое управление}
% **********************************************************************************************



Пусть уравнение состояния динамической системы представляет собой систему стохастических дифференциальных уравнений Ито

\beq{eq:3/2/1}
	d\m{x}_t = \m{a}(\m{x}_t, \m{u}, t)dt + \m{\Sigma}(\m{x}_t, \m{u}, t)d\m{f}_t \text{,}
\eeq

где $\m{x}_t \in \Rv{n}$~--- случайный процесс, являющийся решением уравнения~\ref{eq:3/2/1}; $\m{a} \in \Rv{n}$, $\m{\Sigma} \in \Rv{n}$ представляют собой векторные функции сноса и диффузии соответственно; $\m{f}_t$~--- стандартный винеровский процесс, моделирующий белый шум; $\m{u} \in \Rv{k}$~--- управляющее воздействие, которое выбирается из условия минимизации функционала вида

\beq{eq:3/2/2}
	\funcF(\m{x}, \m{u}, s) = \E \Biggl\{ Q(\m{x}_\tau, \tau) + \int\limits_s^\tau F(\m{x}_t, \m{u}, t)\,dt~\Biggm|~ \m{x}_s = \m{x} \Biggr\} \text{,}
\eeq

где $Q$, $F$~--- известные функции, а $\tau$, например, постоянный момент времени.

Известны различные типы управляющих воздействий, которые могут использовать или не использовать информацию о векторе состояния системы в предшествующие моменты времени (некоторые из них были рассмотрены в предыдущем разделе). Рассмотрим управление $\m{u}(\m{x}_t, t)$ марковского типа. Тогда для него минимальная по $u$ величина функционала $\funcF$, которую обозначим через $\funcH(\m{x}, s)$, удовлетворяет уравнению Гамильтона-Якоби-Беллмана вида

\beqarr
	\lbl{eq:3/2/3}
		\min\limits_{\m{u} \in \Rv{k}} \bigl\{ F(\m{x}, \m{u}, s) + L_u \funcH(\m{x}, s) \bigr\} = 0 \text{;} \\
	\lbl{eq:3/2/4}
		\funcH(\m{x}, T) = Q(\m{x}, T) \text{,}
\eeqarr

где предполагается, что $\tau = T$, а оператор $L_u\argEmpty$ имеет вид

\beq{eq:3/2/5}
\begin{split}
	L_u\argEmpty = \partdiff{\argEmpty}{s} &+ \sum\limits_{i=1}^n\m{a}\idxB{i}(\m{x}, \m{u}, s)\partdiff{\argEmpty}{\m{x}\idxB{i}} + \\
	&+ \frac{1}{2} \sum\limits_{i=1}^n \sum\limits_{j=1}^n \m{D}\idxB{ij}(\m{x}, \m{u}, s)\partdiffsecond{\argEmpty}{\m{x}\idxB{i}}{\m{x}\idxB{j}} \text{,}
\end{split}
\eeq

где $\m{D} = \m{\Sigma}\m{\Sigma}^T$.

Отметим, что приведенные выше условия являются необходимыми условиями, которым удовлетворяет минимум функционала, если он существует.

\br

Рассмотрим линейно-квадратичную проблему стохастического управления, то есть проблему нахождения управления $\m{u}(\m{x}_t, t)$ марковского типа, которое доставляет минимум следующему функционалу:

\beq{eq:3/2/6}
	\funcF(\m{x}, \m{u}, s) = \E \Biggl\{ \m{x}_T\m{R}\m{x}_T^T + \int\limits_s^T \bigr( \m{x}_t^T\m{C}(t)\m{x}_t + \m{u}^T\m{P}(t)\m{u} \bigl)\,dt~\Biggm|~ \m{x}_s = \m{x} \Biggr\} \text{,}
\eeq

где $\m{C}(t), \m{R} \in \Rs{n}{n}$, $\m{P}(t) \in \Rs{k}{k}$, а $\m{x}_t \in \Rv{n}$~--- решение следующей системы стохастических дифференциальных уравнений.

\beq{eq:3/2/7}
	d\m{x}_t = \bigl( \m{A}(t)\m{x}_t + \m{B}(t)\m{u}(\m{x}_t, t) \bigr)dt + \m{\Sigma}(t)d\m{w}_t \text{,}
\eeq

где $\m{A}(t) \in \Rs{n}{n}$, $\m{B}(t) \in \Rs{n}{k}$, $\m{\Sigma} \in \Rs{n}{m}$, а $\m{w}_t \in \Rv{m}$~--- стандартный винеровский процесс с независимыми компонентами $\m{w}\idxB{i}_t$ ($i = 1, 2, \ldots, m$). Кроме этого, предполагается, что все зависящие от $t$ матрицы непрерывны. $\m{C}(t)$ и $\m{R}$~--- симметричные неотрицательно определенные матрицы, а матрица $\m{P}(t)$~--- симметричная положительно определенная.

Запишем уравнение Гамильтона-Якоби-Беллмана для данной задачи:

\beq{eq:3/2/8}
\begin{split}
	\partdiff{\funcH}{s} + \m{x}^T\m{C}(s)\m{x} &+ \sum\limits_{j=1}^n \bigl( \m{A}(s)\m{x} \bigr)\idxB{j}\partdiff{\funcH}{\m{x}\idxB{j}} + \frac{1}{2} \partdiffsecond{\funcH}{\m{x}\idxB{i}}{\m{x}\idxB{j}} + \\
	&+ \min\limits_{\m{u} \in \Rv{k}} \Bigl\{ \m{u}^T\m{P}(t)\m{u} + \sum\limits_{j=1}^n \bigl( \m{B}(s)\m{u} \bigr)\idxB{j}\partdiff{\funcH}{\m{x}\idxB{j}} \Bigr\} = 0 \text{,}
\end{split}
\eeq

где $\m{D} = \m{\Sigma}(t)\m{\Sigma}^T(t)$, а граничное условие имеет вид:

\beq{eq:3/2/9}
	\funcH(T, \m{x}) = \m{x}^T\m{R}\m{x} \text{.}
\eeq

В\cite{KLOEDEN}, в частности, показано, что оптимальное управление и минимальное значение функционала $\funcF(\m{x}, \m{u}, s)$ имеют в этой задаче следующий вид:

\beqarr
	\lbl{eq:3/2/10}
		\m{u}(\m{x}, s) = \m{P}^{-1}(s)\m{B}(s)\m{x} \text{;} \\
	\lbl{eq:3/2/11}
		\funcH(\m{x}, s) = \m{x}^T\m{S}\m{x} + \int\limits_s^T \tr \bigl\{ \m{\Sigma}(t)\m{\Sigma}^T(t)\m{S}(t) \bigr\}\,dt \text{,}
\eeqarr

где $\tr$~--- след\footnote{ Сумма диагональных элементов. } матрицы; $\m{S}(s) \in \Rs{n}{n}$~--- симметричная положительно определенная и непрерывно дифференцируемая матрица, которая удовлетворяет следующему матричному уравнению Риккати:

\beq{eq:3/2/12}
	\genericdiff{S(s)}{s} = -\m{A}(S)^T\m{S}(s) - \m{S}(s)\m{A}(s) + \m{S}(s)\m{B}(s)\m{P}^{-1}(s)\m{B}(s)\m{S}(s) - \m{C}(s)
\eeq

с начальным условием

\beq{eq:3/2/13}
	\m{S}(T) = \m{R} \text{.}
\eeq

Таким образом, если удается решить (хотя бы численно) уравнение Риккати~\ref{eq:3/2/12}, то по формулам~\ref{eq:3/2/10} и~\ref{eq:3/2/11} можно найти оптимальное управление и минимальное значение функционала.

\br

Предположим, что измерение вектора $\m{x}_t$ невозможно, но допустимо измерять наблюдаемый процесс $\m{y}_t \in \Rv{k}$ ($k = 1, 2, \ldots, n$), который удовлетворяет следующей системе дифференциальных уравнений Ито:

\beq{eq:3/2/14}
	d\m{y}_t = \m{H}\m{x}_tdt + \Xi d\m{f}_t
\eeq

с начальным условием $\m{y}_0 = \m{y}(0) = 0$.

В системе~\ref{eq:3/2/14} $\m{H} \in \Rs{k}{n}$, $\Xi \in \Rs{k}{l}$~--- числовые матрицы; $\m{f}_t \in \Rv{l}$~--- стандартный винеровский процесс с независимыми компонентами $\m{f}\idxB{i}_t$ ($i = 1, 2, \ldots, l$), который не зависит от процесса $\m{w}_t$. В этом случае, с помощью фильтра Калмана-Бьюси можно посмтроить оценку $\hat{\m{x}}_t$ и использовать ее в оптимальном управлении $\m{u}(\hat{\m{x}}_t, t)$ для линейно-квадратичной задачи.

Такие исследования остаются за пределами рассмотрения этой работы, подробно они излагаются в соответствующей литературе\cite{KLOEDEN}.

\br

Отметим, что полученные результаты практически полностью, вплоть до переобозначений, совпадают с результатами~\ref{eq:2/2/7},~\ref{eq:2/2/8} и~\vref{eq:2/2/9}, что позволяет говорить о \emph{принципе полной эквивалентности в стохастическом управлении}. Иными словами, построение регулятора, синтезирующего оптимальное управление с обратной связью в случае процессов Ито ничем не отличается от классической модели.



% **********************************************************************************************
\section{Стохастическая устойчивость}
% **********************************************************************************************



Рассмотрим стохастическое дифференциальное уравнение Ито вида~\vref{eq:3/2/1}. Будем считать, что $\m{x}_t \equiv \m{0}$\footnote{ Здесь и далее полагаем <<$\m{0}$>> вектором соответствующей размерности, состоящим из нулей. } является невозмущенным движением динамической системы, а возмущенное движение описывается~\ref{eq:3/2/1}. В соответствии с этим положим

\beq{eq:3/3/1}
    \eqsystem{
		\m{a}(\m{0}, t) \equiv 0 \text{,} \\
		\m{\Sigma}(\m{0}, t) \equiv 0    
    }
\eeq

при $t \geqslant 0$.

\br

Приведем некоторые определения основных типов устойчивости для стохастических систем.

\begin{df}\lbl{df:10}
	Невозмущенное движение системы $\m{x}_t = \m{0}$ будем называть устойчивым по вероятности или стохастически устойчивым при $t \geqslant t_0$, если для любых $\eps > 0$ и $t_0 \geqslant 0$

	\beqn
		\lim\limits_{\m{x}_0 \to \m{0}} \prob \biggl( \sup\limits_{t \geqslant t_0} \bigl| \m{x}^{\m{x}_0, t_0}_t \bigr| \geqslant \eps \biggr) = 0 \text{,}
	\eeqn

	где $\m{x}^{\m{x}_0, t_0}_t$~--- такое решение уравнения~\ref{eq:3/2/1}, что $\m{x}_{t_0} = \m{x}_0$.
\end{df}

\begin{df}\lbl{df:11}
	Невозмущенное движение системы $\m{x}_t \equiv \m{0}$ называется асимптотически устойчивым по вероятности или асимптотически стохастически устойчивым, если оно стохастически устойчиво и, кроме того,
	
	\beqn
		\lim\limits_{\m{x}_0 \to \m{0}} \prob \biggl( \lim\limits_{t \to \infty} \bigl| \m{x}^{\m{x}_0, t_0}_t \bigr| = 0 \biggr) = 1 \text{.}
	\eeqn
\end{df}

Приведенные определения являются обобщением классических определений устойчивости в малом по Ляпунову на случай стохастических систем.

\begin{df}\lbl{df:12}
	Невозмущенное движение системы $\m{x}_t \equiv \m{0}$ называется асимптотически стохастически устойчивым в целом, если оно стохастически устойчио и, кроме того, для всех $|\m{x}_0| < \infty$
	
	\beqn
		\prob \biggl( \lim\limits_{t \to \infty} \bigl| \m{x}^{\m{x}_0, t_0}_t \bigr| = 0 \biggr) = 1 \text{.}
	\eeqn
\end{df}

Приведем также некоторые определения устойчивости, которые связаны с моментными характеристиками процесса $\m{x}_t$.

\begin{df}\lbl{df:13}
	Невозмущенное движение системы $\m{x}_t \equiv \m{0}$ называется устойчивым в среднем степени $p$ ($p > 0$) или $p$-устойчивым, если для любых $\eps > 0$ и $t_0 \geqslant 0$ найдется $\delta = \delta(t_0, \eps) > 0$ такое, что
	
	\beqn
		\E \Bigl\{ \bigl| \m{x}^{\m{x}_0, t_0}_t \bigr|^p \Bigr\} < \eps \text{ для всех } t \geqslant t_0 \text{ и } |\m{x}_0| < \delta \text{.}
	\eeqn
\end{df}

\begin{df}\lbl{df:14}
	Невозмущенное движение системы $\m{x}_t \equiv \m{0}$ называется асимптотически устойчивым в среднем степени $p$ ($p > 0$) или асимптотически $p$-устойчивым, если оно $p$-устойчиво и существует $\delta_0 = \delta_0(t_0) > 0$, такое, что
	
	\beqn
		\lim\limits_{t \to \infty} \E \Bigl\{ \bigl| \m{x}^{\m{x}_0, t_0}_t \bigr|^p \Bigr\} = 0 \text{ для всех } |\m{x}_0| < \delta_0 \text{.}
	\eeqn
\end{df}

\begin{df}\lbl{df:15}
	Невозмущенное движение системы $\m{x}_t \equiv \m{0}$ называется экспоненциально $p$-устойчивым, если существуют постоянные $\gamma > 0$ и $\alpha > 0$, такие, что
	
	\beqn
		\E \Bigl\{ \bigl| \m{x}^{\m{x}_0, t_0}_t \bigr|^p \Bigr\} < \gamma|\m{x}_0|^pe^{-\alpha(t-t_0)} \text{ для всех } t \geqslant t_0 \text{.}
	\eeqn
\end{df}

Следует заметить, что существуют и другие определения стохастической устойчивости\cite{KUSHNER}.

\br

Как известно, одним из методов исследования устойчивости детерменированых систем является метод функций Ляпунова. Достоинство этого метода заключается в том, что он позволяет (в том случае, если соответствующая функция Ляпунова найдена) делать заключения о характере поведения решения дифференциального уравнения, не требуя при этом знания точного решения этого дифференциального уравнения.

Рассмотрим суть метода стохастических функций Ляпунова. Рассмотрим функцию $\funcV(\m{x}, t): \Rv{n} \times \{t > 0\} \to \Rv{1}$, которая при всех $(\m{x}, t) \in \setU \times \{t > 0\}$ ($\setU \subset \Rv{n}$, $\m{0} \in \setU$) является непрерывной и положительно определенной. Пусть также $\funcV(\m{x}, t)$ дважды непрерывно дифференцируема по $\m{x}$ и один раз по $t$, за исключением может быть множества $\{ \m{x} = \m{0} \}$. Кроме этого, предположим, что $\funcV(\m{0}, t) = 0$.

Далее через $\setC_2^0(\setU)$ будем обозначать множество функций дважды непрерывно дифференцируемых по $\m{x} \in \setU$ и один раз по $t$, за исключением множества $\{ \m{x} = \m{0} \}$. В силу формулы Ито имеем

\beq{eq:3/3/2}
\begin{split}
	\funcV(\m{x}_t, t) - \funcV(\m{x}_{t_0}, t_0) &= \int\limits_{t_0}^t L\bigl\{ \funcV(\m{x}_\tau, \tau) \bigr\}\,d\tau + \\
	&+ \sum\limits_{i=1}^m \int\limits_{t_0}^t \m{G}\idxB{i}_0\bigl\{ \funcV(\m{x}_\tau, \tau) \bigr\}\,d\m{f}\idxB{i}_\tau \text{ с вероятностью 1,}
\end{split}
\eeq

где

\beq{eq:3/3/3}
\begin{split}
	L\argEmpty = \partdiff{\argEmpty}{t} &+ \sum\limits_{i=1}^n \m{a}\idxB{i}(\m{x}, t)\partdiff{\argEmpty}{\m{x}\idxB{i}} + \\
	&+ \frac{1}{2}\sum_{j=1}^m\underset{l, i = 1}{\overset{n}{\sum\sum}} \m{\Sigma}\idxB{lj}(\m{x}, t) \m{\Sigma}\idxB{ij}(\m{x}, t) \partdiffsecond{\argEmpty}{\m{x}\idxB{l}}{\m{x}\idxB{i}}
\end{split}
\eeq

и при $i = 1, 2, \ldots, m$

\beq{eq:3/3/4}
	G\idxB{i}\argEmpty = \sum_{j=1}^n \m{\Sigma}\idxB{ji}(\m{x}, t) \partdiff{\argEmpty}{\m{x}\idxB{j}} \text{.}
\eeq

Как можно видеть, в стохастическом случае роль первой производной $\funcV(\m{x}, t)$ по $\m{x}$ в силу уравнения движения играет $L\bigl\{ \funcV(\m{x}_s, s) \bigr\}$. Будем считать, что при $\m{x} \neq \m{0}$ и $t \geqslant 0$ выполняется неравенство

\beq{eq:3/3/5}
	L\bigl\{ \funcV(\m{x}_s, s) \bigr\} \leqslant 0 \text{.}
\eeq

Функцию $\funcV(\m{x}, t)$, обладающую перечисленными выше свойствами, будем называть \emph{стохастической функцией Ляпунова}.

На стохастическую функцию Ляпунова могут также накладываться и некоторые дополнительные ограничения. Далее, используя~\ref{eq:3/3/2} и~\ref{eq:3/3/5} и свойства интеграла Ито, получаем

\beq{eq:3/3/6}
	\E \bigl\{ \funcV(\m{x}_t, t) \bigm| \funcF_{t_0} \bigr\} \leqslant \funcV(\m{x}_{t_0}, t_0) \text{ с вероятностью 1.}
\eeq

Таким образом, $\funcV(\m{x}_t, t)$~--- супермартингал. Если $\m{y}_s$ ($s \geqslant 0$)~--- мартингал и $\E \bigl\{ |\m{y}_s|^p \bigr\} < \infty$, то справедливо хорошо известное неравенство:

\beq{eq:3/3/7}
	\forall \alpha:~\prob\Biggl(\bigl\{ \omega \in \Omega : \sup\limits_{0 \leqslant s \leqslant t} |\m{y}(s, \omega)| \geqslant \alpha \bigr\}\Biggr) \leqslant \frac{1}{\alpha^p} \E \bigl\{ |\m{y}_t|^p \bigr\}
\eeq

и, следовательно, для всех $T > t_0$ и $\delta > 0$

\beq{eq:3/3/8}
	\prob\Biggl( \sup\limits_{t_0 \leqslant t \leqslant T} \funcV(\m{x}_t, t) \geqslant \delta \Biggr) \leqslant \frac{1}{\delta} \funcV(\m{x}_{t_0}, t_0) \text{.}
\eeq

Из последнего неравенства, непрерывности $\funcV(\m{x}, t)$ и того, что $\funcV(\m{0}, t) = 0$, получаем:

\beq{eq:3/3/9}
	\lim\limits_{\m{x}_0 \to \m{0}} \prob\Biggl( \sup\limits_{t_0 \leqslant t \leqslant T} \funcV(\m{x}_t, t) \geqslant \delta \Biggr) = 0 \text{.}
\eeq

Далее, если предположить, что при некоторых постоянных $p$, $k_1$ и $k_2$

\beq{eq:3/3/10}
	k_1 |\m{x}|^p \leqslant \funcV(\m{x}, t) \leqslant k_2 |\m{x}|^p \text{,}
\eeq

то согласно~\ref{eq:3/3/9} получим, что невозмущенное движение системы $\m{x}_t~\equiv~\m{0}$ стохастически устойчиво.

\br

Приведем еще ряд фактов\cite{HASMINSKI}, касающихся стохастической устойчивости

\bteo{teo:5}
	Пусть в области $\{ t > 0 \} \times \setU$, включающей множество $\{ \m{x} \equiv \m{0} \}$, существует непрерывная, положительно определенная функция $\funcV(\m{x}, t) \in \setC_2^0(\setU)$ и пусть при $t > 0$
	
		\beqn
			\limsup\limits_{\m{x} \to \m{0}} \funcV(\m{x}, t) = 0 \text{,}
		\eeqn
		
		а $L\bigl\{ \funcV(\m{x}, t) \bigr\} < 0$ в этой области. Тогда невозмущенное движение $\m{x}_t \equiv \m{0}$ асимптотически стохастически устойчиво.
\eteo

\bteo{teo:6}
	Пусть для всех $\m{x} \in \Rv{n}$ существует положительно определенная функция $\funcV(\m{x}, t) \in \setC_2^0(\Rv{n})$, для которой при $t > 0$
	
	\beqn
		\eqsystem{
			\limsup\limits_{\m{x} \to \m{0}} \funcV(\m{x}, t) = 0 \text{,} \\
			\liminf\limits_{\m{x} \to \m{\infty}} \funcV(\m{x}, t) = \infty \text{,} \\
			L\bigl\{ \funcV(\m{x}, t) \bigr\} < 0 \text{.}
		}
	\eeqn
	
	Тогда невозмущенное движение $\m{x}_t \equiv \m{0}$ асимптотически стохастически устойчиво в целом. 
\eteo

\bteo{teo:7}
	Для экспоненциальной $p$-устойчивости невозмущенного движения $\m{x}_t \equiv \m{0}$ при $t > 0$ достаточно, чтобы при всех $\m{x} \in \Rv{n}$ существовала функция $\funcV(\m{x}, t) \in \setC_2^0(\Rv{n})$, удовлетворяющая при некоторых положительных постоянных $k_1$, $k_2$ и  $k_3$ неравенствам
	
	\beqn
		\eqsystem{
			k_1 |\m{x}|^p \leqslant \funcV(\m{x}, t) \leqslant k_2 |\m{x}|^p \text{,} \\
			L\bigl\{ \funcV(\m{x}, t) \bigr\} < -k_3 |\m{x}|^p \text{.}
		}
	\eeqn
\eteo

Другой метод детерминистской теории устойчивости заключается в линеаризации дифференциального уравнения в окрестности невозмущенного движения и анализе устойчивости нулевого решения линеаризованного дифференциального уравнения. Известно\cite{ARNOLD}, что в ряде случаев устойчивость нулевого решения линеаризованной системы влечет устойчивость невозмущенного решения соответствующей нелинейной системы. Особую роль при анализе устойчивости с помощью этого подхода играют так называемые ляпуновские экспоненты или ляпуновские показатели. Оказывается, что метод линеаризации может быть обобщен и на случай стохастических систем\cite{AUSLENDER}.


% пример типовой задачи
% coding: utf-8
% --------------------------------------------------------------------------------------------------
% "Синтез оптимального стохастического управления", 2011 год
% --------------------------------------------------------------------------------------------------



\chapter{Пример решения типовой задачи}
% ==============================================================================================
\renewcommand{\optU}{  \optimum{\m{u}}  } % оптимальное управление
\renewcommand{\funcF}{ \calf{F}         } % множество F

\renewcommand{\optF}{  \optimum{\funcF} } % оптимальный критерий качества
\renewcommand{\funcS}{ \calf{S}         } % второстепенный критерий качества
\renewcommand{\funcI}{ \calf{I}         } % метод сравнения
% ==============================================================================================



% **********************************************************************************************
\section{Постановка задачи}
% **********************************************************************************************



Рассмотрим полученные в разделах 2 и 3 результаты, применим их к конкретной задаче.

Пускай стоит задача вывода на орбиту искусственного спутника. В подобных задачах ракета-носитель должна следовать по номинальной траектории. Но из-за нестационарности параметров, этого не будет. Для коррекции небольших отклонений относительно заданной траектории можно, например, линеаризовать динамику ракеты-носителя относительно заданной траектории. После выбора критерия, можно приступать к расчетам в соответствии с методиками этого раздела.

Динамику линеаризованного контура наведения ракеты можно выразить в следующей форме:

\beq{eq:4/1/1}
	\eqsystem{
		\dot{x}_1 &= x_2 \mbox{,} \\
		\dot{x}_2 &= \frac{k_1}{k_2 - t} x_3 \mbox{,} \\
		\dot{x}_3 &= u \mbox{.}
	}
\eeq

где $0 \leqslant t \leqslant T$ и

\begin{description}
	\item[$x_1$]~--- боковое отклонение от номинальной траектории;
	\item[$x_2$]~--- скорость этого отклонения;
	\item[$x_3$]~--- угол направления вектора тяги.
\end{description}

Зависимость между боковой тягой и боковым ускорением нестационарна и определяется коэффициентом $k_1 / (k_2 - t)$, который учитывает потерю массы при действии тяги. Интегрируемая связь между $x_3$ и $u$ представляет собой линеаризованное уравнение привода.

%\addfigure{fig:1}{x_adaptive_noise_high}{Функции матрицы $\m{P}$.}

Если исходить из предположения о возможности точного измерения $x_1(t)$, $x_2(t)$ и $x_3(t)$, то можно разработать такую систему управления с обратной связью, чтобы обеспечить оптимальное наведение ракеты, удовлетворив при этом заданному показателю качества. Предположим, что используется функционал вида

\beq{eq:4/1/2}
	\funcF = \frac{1}{2} \int\limits_0^T \bigl( \m{x}^T(\tau)\m{Q}(\tau)\m{x}(\tau) + ru^2(t) \bigr)\,d\tau \mbox{.}
\eeq

Систему~\ref{eq:4/1/1} можно представить в обозначениях раздела 2.2, а именно в виде системы $\m{x}(t) = \m{A}(t)\m{x}(t) + \m{B}(t)\m{u}(t)$. Это возможно, поскольку матрицы $\m{A}$ и $\m{B}$ выражаются как

\beqarr
	\lbl{eq:4/1/3}
		\m{A}(t) = \matr{
			0 & 1 & 0 \\
			0 & 0 & \frac{k_1}{k_2-t} \\
			0 & 0 & 0		
		} \mbox{;} \\
	\lbl{eq:4/1/4}
		\m{B}(t) = \m{b} = \matr{
			0 \\ 0 \\ 1		
		} \mbox{.}
\eeqarr

Иными словами, задача выражается как задача из раздела 2.2 с критерием качества~\vref{eq:2/2/11}. Как было показано, для ее решения, необходимо получить матрицу $\m{P}(t)$ с граничным условием $\m{P}(T) = \m{0}$. Как было показано ранее, оптимальное уравнение в данном случае равно

\beq{eq:4/1/5}
	\optU(t) = -\frac{1}{r}\m{B}(t)^T\m{P}(t)\m{x}(t) = -\frac{1}{r}\m{b}^T\m{P}(t)\m{x}(t) \mbox{.}
\eeq

С другой стороны, так как

\beq{eq:4/1/6}
	\m{P}(t) = \matr{
		p_{11}(t) & p_{21}(t) & p_{31}(t) \\
		p_{12}(t) & p_{22}(t) & p_{32}(t) \\
		p_{13}(t) & p_{23}(t) & p_{33}(t) \\
	} \mbox{,}
\eeq

то учитывая~\ref{eq:4/1/4} и~\ref{eq:4/1/5}, получаем, что

\beq{eq:4/1/7}
	\optU(t) = -\frac{1}{r} \matr{ p_{13}(t) & p_{23}(t) & p_{33}(t) } \m{x}(t) \mbox{.}
\eeq

\br

Учтем теперь влияние атмосферы на движение ракеты. Скорость ветра $\m{V}$, вообще говоря, является векторной случайной функцией, состоящей из координат точки атмосферы и времени

\beq{eq:4/1/8}
	\m{V} = \m{V}(p, q, z, t) \mbox{,}
\eeq

где $p$, $q$, $z$~--- пространственные координаты\footnote{Традиционное обозначение координат как $x_0$, $y_0$ и $z_0$ в работе не используется, поскольку эти символы уже заняты.}, а $t$~--- время. Иными словами, скорость ветра является векторным случайным полем.

Величина $\m{V}$ может быть представлена в виде суммы среднего вектора $\bar{\m{V}}$ и центрированной составляющей $\m{V}_{\mbox{сл}}$:

\beq{eq:4/1/9}
	\m{V} = \bar{\m{V}} + \m{V}_{\mbox{сл}} (p, q, z, t) \mbox{.}
\eeq

Под $\bar{\m{V}}$ понимается среднее по некоторому интервалу времени значение вектора ветра в рассматриваемой точке. Наличие случайной составляющей $\m{V}_{\mbox{сл}}$ означает турбулентность атмосферы.

Вообще говоря, колебания или вариации скорости ветра в атмосфере имеют очень широкий спектр, начиная от годовых изменений и кончая пульсациями с частотой в десятки и даже сотни герц, вызванных самыми мелкими вихрями вблизи поверхности земли или даже самого летательного аппарата. Однако для задач управления движением весьма низкочастотные вариации не имеют значения\cite{KRASOVSKY}, и усреднение в формуле~\ref{eq:4/1/9} может осуществляться по интервалу, измеряемому минутами.

Если все статистические характеристики случайной функции $\m{V}$, вычисляемые для произвольной точки $(p_0, q_0, z_0)$, не зависят от этой точки, то есть одинаковы для рассматриваемого пространства, то турбулентность называется \emph{однородной}. Некоторые статистические характеристики, такие как корреляционные функции или законы совместного распределения вероятностей, вычисляются для двух или большего числа точек рассматриваемого пространства. Если эти статистические характеристики не меняются при параллельных переносах, вращениях и зеркальных отражениях этой системы, то турбулентность называется \emph{изотропной}. Иными словами, поле турбулентности изотропно, если в нем нельзя выделить никакого преимущественного направления.

Принято полагать, что анизотропность турбулентности атмосферы вызыается преимущественно рельефом местности. Для высот, больших 1000 метров, турбулентность атмосферы считают изотропной. При рассмотрении движения ракеты скорость воздуха можно считать случайной функцией не трех, а только одной пространственной координаты $p$, отсчитываемой вдоль траектории полета:

\beq{eq:4/1/10}
	\m{V}_{\mbox{сл}} = \m{V}_{\mbox{сл}} (p, t) \mbox{.}
\eeq

Объясняется это тем, что на движение ракеты как твердого тела слабо влияют высокочастотные составляющие\cite{KRASOVSKY}, создаваемые вихрями меньших размера самой ракеты. Поэтому скорость ветра в пределах контуров тела приближенно можно считать одинаковой. Для аэроупругого летательного аппарата обычно необходимо учитывать более высокочастотные составляющие воздействий турбулентной атмосферы, следовательно уравнение~\ref{eq:4/1/10} не срабатывает.

Возвращаясь к привычным обозначениям, можно считать, что случайная составляющая зависит лишь от времени $t$ и бокового отклонения от номинальной траектории $x_1$:

\beq{eq:4/1/11}
	\m{V}_{\mbox{сл}} = \m{V}_{\mbox{сл}} \bigl(x_1(t), t \bigr) \mbox{.}
\eeq

Впрочем, при достаточно большой скорости самолета вторым аргументом можно пренебречь, полагая

\beq{eq:4/1/12}
	V_{\mbox{сл}} = \m{V}_{\mbox{сл}} \bigl(x_1(t) \bigr) \mbox{.}
\eeq

Это означает, что распределение ветра в пространстве считается как бы <<замороженным>> и зависимость от времени получается за счет того, что самолет пронизывает это распределение. Данное допущение носит название \emph{гипотезы Тейлора}\cite{MONIN}. Гипотеза Тейлора подвергалась экспериментальной проверке неоднократно, противоречащих результатов получено не было.

Скорость ветра, очевидно, влияет на скорость отклонения ракеты с траектории, следовательно система~\vref{eq:4/1/1} может быть переписана с поправкой на турбулентность атмосферы следующим образом:

\beq{eq:4/1/13}
	\eqsystem{
		\dot{x}_1 &= x_2 \mbox{,} \\
		\dot{x}_2 &= \frac{k_1}{k_2 - t} x_3 + V(x_1) \mbox{,} \\
		\dot{x}_3 &= u \mbox{.}
	}
\eeq

Будем считать, что пространство изотропно, а скорость ракеты достаточна для того, чтобы можно было воспользоваться гипотезой Тейлора. Не нарушая общности, можно полагать, средний по интервалу вектор $\m{V}$ равен нулевому. Иными словами, мы предполагаем такую ситуацию, когда нет перепадов давления, и, следовательно, нет постоянного ветра (либо скорость ветра пренебрежительно мала по сравнению со скоростью ракеты), лишь некоторые атмосферные шумы. Небольшие естественные атмосферные воздействия можно представить в виде белого шума, который моделируется стандартным винеровским процессом, определение которому дано на странице~\pageref{df:6}.

Таким образом, систему~\ref{eq:4/1/13} можно представить в виде следующего процесса Ито:

\beq{eq:4/1/14}
	\eqsystem{
		\dot{x}_1 &= x_2 \mbox{,} \\
		\dot{x}_2 &= \frac{k_1}{k_2 - t} x_3 + \xi(t) \mbox{,} \\
		\dot{x}_3 &= u \mbox{.}
	}
\eeq

Функция $\frac{k_1}{k_2 - t} x_3$ в уравнении для $\dot{x}_2$ будет определять коэффицент сноса, а случайная функция $\xi(t)$~--- коэффицент диффузии.

Задачу~\ref{eq:4/1/14} можно моделировать следующим образом

\beq{eq:4/1/15}
	\eqsystem{
		\dot{x}_1 &= x_2 \mbox{,} \\
		\dot{x}_2 &= \frac{k_1}{k_2 - t} x_3 + \xi_t \mbox{,} \\
		\dot{x}_3 &= u \mbox{.}
	}
\eeq

где $\{\xi_t\}_{t=0}^\infty$~--- последовательность независимых одинаково распределенных случайных величин с распределением Гаусса. Известно\cite{ARKHIPOV}, что в этом случае случайный процесс~\ref{eq:4/1/15} будет слабо сходиться к процессу~\ref{eq:4/1/15}.



% **********************************************************************************************
\section{Исследование и анализ алгоритмов решения поставленной задачи}
% **********************************************************************************************



Рассмотрим систему~\ref{eq:4/1/15} c критерием качества~\ref{eq:4/1/2}.

Для того, чтобы получить матрицу $\m{P}(t)$ требуется решить уравнение~\vref{eq:2/2/7}. %Непосредственно решая его, в силу~\ref{eq:4/1/3} и~\ref{eq:4/1/4} получаем следующую систему нелинейных дифференциальных уравнений с граничными условиями в одной точке:

%\beq{eq:4/2/1}
%	\eqsystem{
%		\dot{p}_{11}(t) &= \frac{\mathstrut 1}{10}p^2_{13}(t) - \frac{5 \cdot 10^{-7}}{(300-t)^2}         \mbox{,} \\
%		\dot{p}_{12}(t) &= \frac{1}{10}p_{13}(t)p_{23}(t) - p_{11}(t)                                     \mbox{,} \\
%		\dot{p}_{13}(t) &= \frac{1}{10}p_{13}(t)p_{33}(t) - \frac{k_1}{k_2-t} p_{12}(t)                   \mbox{,} \\
%		\dot{p}_{21}(t) &= 0                                                                              \mbox{,} \\
%		\dot{p}_{22}(t) &= \frac{1}{10}p^2_{23}(t) - 2p_{12}(t) - \frac{10^{-3}}{(300-t)^2}               \mbox{,} \\
%		\dot{p}_{23}(t) &= \frac{1}{10}p_{23}(t)p_{33}(t) - \frac{k_1}{k_2-t}p_{22}(t) - p_{13}(t)        \mbox{,} \\
%		\dot{p}_{31}(t) &= 0                                                                              \mbox{,} \\
%		\dot{p}_{32}(t) &= 0                                                                              \mbox{,} \\
%		\dot{p}_{33}(t) &= \frac{1}{10}p^2_{33}(t) - 2\frac{k_1}{k_2-t}p_{23}(t) - \frac{10^3}{(300-t)^2} \mbox{,}
%	}
%\eeq

Эту задачу, вследствие нелинейности системы нужно решать численно, двигаясь в обратном времени от $t=T$ с граничным условием $\m{P}(T) = \m{0}$.

Таким образом, можно сформулировать первый алгоритм решения задачи:

\balgo{alg:1}
	\benum
		\item
			Для заданных $T$, $r$, $\m{Q}(t)$ требуется решить матричное дифференциальное уравнение~\ref{eq:2/2/7} в обратном времени с граничным условием $\m{P}(T) = \m{0}$. Таким образом, получим траекторию $\m{P}(t)$, где $0 \leqslant t \leqslant T$;
			
		\item
			Пользуясь формулой~\ref{eq:4/1/7} получаем оптимальное управление $\optU(t)$ на интервале $0 \leqslant t \leqslant T$ в каждый момент времени $t$\footnote{Обратим внимание, что таким образом синтезируется управление с обратной связью, следовательно необходимо знать точные значения каждой переменной состояния из вектора $\m{x}(t)$ в рассматриваемый момент времени.}. 
	\eenum
\ealgo

\br

Рассмотрим теперь случай, когда система стационарна. Как было показано ранее, при $T \to \infty$, уравнение~\vref{eq:2/2/7} сводится к нелинейному алгебраическому уравнению Риккати~\vref{eq:2/2/10}. Остальные действия остаются прежними. Формула~\ref{eq:4/1/7} также преобразуется с учетом этой особенности:

\beq{eq:4/2/2}
	\optU(t) = -\frac{1}{r} \matr{ p_{13} & p_{23} & p_{33} } \m{x}(t) \mbox{.}
\eeq

Следовательно, можно построить новый алгоритм, берущий за основу алгоритм~\ref{alg:1}.

\balgo{alg:2}
	\benum
		\item
			Для заданных $T$, $r$, $\m{Q}$ требуется решить матричное алгебраическое уравнение~\ref{eq:2/2/10}. Таким образом, получим матрицу с постоянными коэффициентами $\m{P}$ для всех $0 \leqslant t \leqslant T$;
			
		\item
			Пользуясь формулой~\ref{eq:4/2/2} получаем оптимальное управление $\optU(t)$ на интервале $0 \leqslant t \leqslant T$ в каждый момент времени $t$.
	\eenum
\ealgo

Опять же, стоит обратить внимание на то, что таким образом синтезируется управление с обратной связью, следовательно необходимо знать точные значения каждой переменной состояния из вектора $\m{x}(t)$ в рассматриваемый момент времени.

\br

Метод Беллмана, неявно лежащий в основе вышеприведенных алгоритмов, в некотором смысле аддитивен. Уравнения Беллмана позволяют найти оптимальное управление не только в целом, но и на конкретном шаге. В основной теореме динамического программирования, доказанной Беллманом утверждается, что управление, построенное таким образом, будет оптимальным. Действительно, одним из следствий уравнения~\vref{eq:2/1/17} является тот факт, что в любой момент времени величина критерия должна быть минимальной. Следовательно, если разбить задачу динамического программирования на несколько частей, прямо следующих друг за другом, причем состояние, в котором система оказалась в конце предыдущей части, будет начальным состоянием последующей, то в результате траектория динамики системы окажется такой, при которой критерий качества будет минимальным, а полученное управление~--- оптимальным.

Из уравнений~\ref{eq:2/1/3} и~\ref{eq:2/1/4}, а также из утверждения~\vref{statement:1} следует, что такое кусочное управление будет эквивалентно оптимальному. Возможна ситуация, когда управление, полученное таким образом, будет отличаться от управления, полученного из уравнений~\ref{eq:2/2/8} и~\vref{eq:2/2/9}, однако ввиду того, что значения критериев равны (оба уравнения доставляют минимум критерию качества~\vref{eq:2/2/1}), оба управления будут являться оптимальными. Иными словами, уравнения Беллмана позволяют получить оптимальное управление, но из них, конечно же, никоим образом не следует единственность минимума критерия качества.

Из этих рассуждений следует, что не имеет большого значения то, каким способом получается оптимальное управление в случае детерминированной системы~\vref{eq:4/1/1}, однако различия в подходах имеют смысл в случае стохастической системы~\vref{eq:4/1/14}. В сильно зашумленных системах, где диффузия оказывает большое влияние на динамику траекторий, дробя процесс вычисления траектории на несколько, можно построить адаптивный алгоритм, который будет учитывать не только параметры системы, но и степень отклонения от намеченной траектории. Иными словами, адаптивный характер алгоритма должен минимизировать влияние тех факторов, которые прямо контролировать невозможно, либо слишком сложно и затратно.

Отметим тот факт, что оба предложеных ранее алгоритма рассчитывают управление сразу на весь временной интервал наблюдений, после чего ракета наводится программно, без учета внешних воздействий. Предлагается рассчитывать управление поэтапно.

\br

Основная идея предлагаемого адаптивного управления заключается в следующем. Разобъем временной промежуток наблюдений $0 \leqslant t \leqslant T$ на $n+1$ равную часть: $0 = t_0 \leqslant t_1 \leqslant t_2 \leqslant \ldots \leqslant t_{n} = T$ ($t_{i+1}-t_i = t_{j+1}-t_j$ для всех $i, j < n$). Дальше, следуя алгоритмам~\ref{alg:1} и~\ref{alg:2}, синтезируем управление на всем интервале наблюдений, однако управление используем лишь на первом временном промежутке $[t_0, t_1)$. В момент времени $t_1$ происходит проверка эффективности проведенного управления с помощью некоторого заранее определенного вторичного критериия качества $\funcS(\m{x}, t_0, t_1)$ по заданному правилу.

Если сравнение показало, что управление не столь эффективно, как бы этого хотелось, то синтезированное ранее управление отбрасывается, и задача решается заново, но уже для промежутка $t_1 \leqslant t \leqslant T$, однако изначальное деление на интервалы сохраняется.

Если же сравнение управления показало его эффективность на первом интервале, то используется синтезированное ранее управление. В обоих случаях, в конце промежутка $[t_1, t_2)$, в момент времени $t_2$ происходит аналогичная проверка вторичного критерия качества $\funcS(\m{x}, t_1, t_2)$.

Такие проверки проводятся вплоть до самого конца временного интервала наблюдений, за исключением лишь момента времени $t_n = T$.

\br

Как можно видеть, идея адаптивного алгоритма активно использует понятие <<доверительного интервала>>. Это означает, что управление не синтезируется от начала до конца, а потом применяется без каких-либо корректировок. Иными словами, ракетой не \emph{управляют}, а \emph{пытаются управлять}, и если результат далек от желаемого, вносятся определенные исправления. Эти корректировки производятся по прошествию определенного времени в том случае, если они имеют смысл: синтезированному управлению более не доверяется, строится новое, принимая за начальное состояние системы то, в котором она оказалась к текущему моменту времени. Если большого смысла в корректировках нет, и отклонения от желаемой траектории укладываются в заданные рамки, то управление считается легитимным, и ракета управляется посредством него следующий интервал времени.

Как можно видеть, такой подход не является композитным. На каждом интервале времени используется управление, полученное с помощью предыдущих двух алгоритмов. Адаптивный алгоритм определяет то, когда требуется пересчет матрицы $\m{P}(t)$, являющейся ключевым элементом синтеза, что не заложено в других методах. Иными словами, подобный подход инкапсулирует предыдущие алгоритмы.

Подобный подход полезен, когда вычисление оптимального управления на длинном интервале отнимает множество ресурсов, в том числе и временных, либо когда система сильно зашумлена. Таким образом, можно рассчитывать данные на короткие интервалы времени и корректировать поведение в соответствии желаемым результатом.

Адаптивный подход оперирует понятием вторичных критериев качества $\funcS(\emptyarg, \emptyarg, \emptyarg)$. Эти критерии должны быть тесно связаны с заданным критерием качества задачи. Они требуются в том случае, когда вычисление основного критерия отнимает слишком много времени, либо сопряжено с иными трудностями. Вторичные критерии качества $\funcS(\emptyarg, \emptyarg, \emptyarg)$ можно полагать как грубые аппроксимации самого критерия качества $\funcF(\emptyarg)$, можно выбирать некоторые эвристические характеристики,~--- все зависит от конкретного случая. От выбора $\funcS(\emptyarg, \emptyarg, \emptyarg)$ будет напрямую зависеть частота пересчета оптимальных управлений. Например, его можно выбрать следующим образом:

\beq{eq:4/2/3}
	\funcS(\m{x}, t_i, t_{i+1}) \eqdef \Biggl| \int\limits_{t_i}^{t_{i+1}} \bigl( \m{x}^T(\tau)\m{Q}(\tau)\m{x}(\tau) + ru(t)^2 \bigr)\,d\tau - 2\optF(\m{x}, t_i, t_{i+1}) \Biggr| \mbox{.}
\eeq

Иными словами, уравнение~\ref{eq:4/2/3} определяет абсолютную разницу между реальной динамикой зашумленной системы и ее же, но в отсутсвии каких-либо случайных внешних воздействий. Очевидно, чем больше эта разница, тем скорее необходимо корректировать управление. В силу уравнения~\vref{eq:4/1/2} получаем, что $\funcS(\m{x}, t_i, t_{i+1}) \geqslant 0$\footnote{Отметим, что хоть траектории могут отличаться, но в случае, когда $\funcS(\m{x}, t_i, t_{i+1}) = 0$, это не является существенным обстоятельством, поскольку значение критерия качества на этом интервале совпадает с оптимальным значением.}. В случае же, когда воздействий не происходит, реальная траектория не будет отличаться от оптимальной, поэтому $\funcS(\m{x}, t_i, t_{i+1}) = 0$.

Таким образом, если ввести некоторое пороговое значение $\eps > 0$, то можно определить ситуацию, когда управлению <<доверяют>>, а когда нет: в случае, если $\funcS(\m{x}, t_i, t_{i+1}) \leqslant \eps$, управлению доверяют, в противном случае требуется пересчет.

Идею вышеприведенного адаптивного метода можно выразить следующим алгоритмом:

\balgo{alg:3}
	\benum
		\item
			Делим временной интервал наблюдений $[0, T]$ на $n+1$ равную часть: $0 = t_0 \leqslant t_1 \leqslant t_2 \leqslant \ldots \leqslant t_{n} = T$, где $t_{i+1}-t_i = t_{j+1}-t_j$ для всех $i, j < n$. Интервалы типа $[t_i, t_{i+1})$ ($i<n$) назовем подынтервалами;
		
		\item
			Строим оптимальное управление $\optU$ на всем оставшемся интервале наблюдений с помощью алгоритмов~\ref{alg:1} или~\vref{alg:2};
		
		\item
			Наводим ракету с помощью управления $\optU$ на протяжении одного подынтервала времени;
		
		\item
			На открытой границе подынтервала производим расчет вторичного критерия $\funcS(\m{x}, t_i, t_{i+1})$ (например, заданного уравнением~\ref{eq:4/2/3});
			
		\item
			Если $\funcS(\m{x}, t_i, t_{i+1}) \leqslant \eps$ ($\eps > 0$), то переходим на шаг 3, в противном случае переходим на шаг 2.

	\eenum
\ealgo

\br

Покажем, что в случае незашумленной системы управление, полученное с помощью алгоритма~\ref{alg:3} будет оптимальным.

\begin{teo}\lbl{teo:5}
    В случае, если динамическая система описывается уравнениями~\ref{eq:4/1/1} и \vref{eq:4/1/2}, а также вторичным критерием качества~\vref{eq:4/2/3}, алгоритм~\ref{alg:3} синтезирует оптимальное управление $\optU$.
\end{teo}

\strong{Доказательство.}

Рассмотрим критерий качества~\ref{eq:4/2/3}. Отметим, что интеграл в этом уравнении совпадает с интегральным критерием качества~\ref{eq:4/1/2}, если бы задача рассматривалась на промежутке времени $[t_i, t_{i+1}]$. В силу свойства аддитивности интегралов, уравнение~\ref{eq:4/1/2} можно представить следующим образом

\beq{eq:4/2/4}
\begin{split}
	\funcF = \frac{1}{2} \Biggl(& \int\limits_{t_0 = 0}^{t_1} \bigl( \m{x}^T(\tau)\m{Q}(\tau)\m{x}(\tau) + ru^2(t) \bigr)\,d\tau + \\
	&+ \int\limits_{t_1}^{t_2} \bigl( \m{x}^T(\tau)\m{Q}(\tau)\m{x}(\tau) + ru^2(t) \bigr)\,d\tau + \\
	&+ \cdots + \\
	&+ \int\limits_{t_{n-1}}^{t_n = T} \bigl( \m{x}^T(\tau)\m{Q}(\tau)\m{x}(\tau) + ru^2(t) \bigr)\,d\tau \Biggr) \mbox{.}
\end{split}
\eeq

При использовании оптимального управления $\funcF \equiv \optF$, следовательно уравнение~\ref{eq:4/2/4} можно представить в виде

\beq{eq:4/2/5}
\begin{split}
	\funcF \equiv \optF &= \frac{1}{2} \Biggl( \int\limits_{t_0 = 0}^{t_1} \bigl( \m{x}^T(\tau)\m{Q}(\tau)\m{x}(\tau) + r\optimum{u}^2(t) \bigr)\,d\tau + \\
	&+ \int\limits_{t_1}^{t_2} \bigl( \m{x}^T(\tau)\m{Q}(\tau)\m{x}(\tau) + r\optimum{u}^2(t) \bigr)\,d\tau + \\
	&+ \cdots + \\
	&+ \int\limits_{t_{n-1}}^{t_n = T} \bigl( \m{x}^T(\tau)\m{Q}(\tau)\m{x}(\tau) + r\optimum{u}^2(t) \bigr)\,d\tau \Biggr) = \\
	& = \frac{1}{2} \sum\limits_{i=0}^{n-1} \optF(\m{x}, t_i, t_{i+1}) = \sum\limits_{i=0}^{n-1} \frac{1}{2} \optF(\m{x}, t_i, t_{i+1}).
\end{split}
\eeq

Это справедливо в силу утверждения~\vref{statement:1}.

Таким образом, критерий качества~\ref{eq:4/1/2} является аддитивным, на каждом подынтервале времени $[t_i, t_{i+1})$ его минимальное значение при использовании оптимального управления, полученного из алгоритмов~\ref{alg:1} или~\ref{alg:2}, равно $\optF(\m{x}, t_i, t_{i+1})$.

Но из этого следует, что вторичный критерий качества, определяемый уравнением~\ref{eq:4/2/3} будет всегда равняться нулю, а значит, что $\funcS(\m{x}, t_i, t_{i+1}) < \eps$ для любого наперед заданного $\eps > 0$. Следовательно, 4 шаг алгоритма~\ref{alg:3} будет приводить к тому, что управление никогда не будет пересчитываться, а значит оно будет ровно тем же, что и при использовании алгоритмов~\ref{alg:1} или~\ref{alg:2}, а, следовательно, оптимальным. \qed

Стоит отметить, что алгоритм~\ref{alg:3} является в некотором смысле эвристическим, поскольку эффективность подобных перерасчетов в общем случае крайне сложно изучить. С другой стороны, в результате большого количества численных экспериментов с наиболее типовыми функциями, не удалось найти систему, когда такой подход давал бы худшие результаты, нежели алгоритмы~\ref{alg:1} и~\ref{alg:2}, даже наоборот: в системах с большими шумами он, с поправкой на погрешность численных методов, дает меньшее отклонение от значения оптимального критерия качества детерменированного случая.

Также следует обратить внимание на то, что алгоритм~\ref{alg:3} с вторичным критерием качества~\ref{eq:4/2/3} имеет гораздо более высокие требования к вычислительным устройствам, поскольку помимо расчета оптимального управления, требуется всякий раз рассчитывать значения двух интегралов. С другой стороны, эти задачи независимы друг от друга и прекрасно могут быть распараллелены, что при наличии двух и более процессоров приводит к сравнимому быстродействию. Впрочем, следует отметить и тот факт, что алгоритм 3 с таким вторичным критерием качества использует почти в два раза большее количество оперативной памяти для своей работы. Впрочем, в случае системы со слабой зашумленностью (или даже в полностью детерменированном случае), никакого существенного изменения в требованиях к ее количеству не происходит.

\br

Обратим внимание на то, что модфицирование алгоритма~\vref{alg:3} с целью разного рода оптимизаций, может привести к неожиданным нежелательным результатам. Рассмотрим следующую модификацию этого алгоритма:

\balgo{alg:4}
	\benum
		\item
			Делим временной интервал наблюдений $[0, T]$ на $n+1$ равную часть: $0 = t_0 \leqslant t_1 \leqslant t_2 \leqslant \ldots \leqslant t_{n} = T$, где $t_{i+1}-t_i = t_{j+1}-t_j$ для всех $i, j < n$. Интервалы типа $[t_i, t_{i+1})$ ($i<n$) назовем подынтервалами. Оптимальное управление строится поподынтервально;
		
		\item
			Строим оптимальное управление $\optU$ только лишь на следующем подынтервале с помощью алгоритмов~\ref{alg:1} или~\vref{alg:2};
		
		\item
			Наводим ракету с помощью управления $\optU$ на протяжении одного подынтервала времени;
		
		\item
			На открытой границе подынтервала производим расчет вторичного критерия $\funcS(\m{x}, t_i, t_{i+1})$ (например, заданного уравнением~\ref{eq:4/2/3});
			
		\item
			Если $\funcS(\m{x}, t_i, t_{i+1}) > \eps$ ($\eps > 0$), то сразу переходим на шаг 2. В противном случае, также переходим на шаг 2, но в качестве следующего подынтервала рассматриваем интервал удвоенной длины, если оставшихся подынтервалов достаточно (если нет~--- берется оставшийся временной промежуток). Иными следующим подынтервалом будет не $[t_{i+1}, t_{i+2})$, а $[t_{i+1}, t_{i+3})$.
	\eenum
\ealgo

Для алгоритма~\ref{alg:4} можно сформулировать теорему, аналогичную теореме~\ref{teo:5}.

\begin{teo}\lbl{teo:6}
    В случае, если динамическая система описывается уравнениями~\ref{eq:4/1/1} и \vref{eq:4/1/2}, а также вторичным критерием качества~\vref{eq:4/2/3}, алгоритм~\ref{alg:4} синтезирует оптимальное управление $\optU$.
\end{teo}

Доказательство теоремы совершенно аналогично доказательству теоремы~\ref{teo:5}.

Однако, алгоритм~\ref{alg:4} совершенно не подходит к задаче управления ракетой в случае даже малого влияния шума. Более подробно этот вопрос будет исследован в разделе 4.4.



% **********************************************************************************************
\section{Программная реализация решения}
% **********************************************************************************************



% **********************************************************************************************
\section{Результаты решения}
% **********************************************************************************************


% заключение
% coding: utf-8
% --------------------------------------------------------------------------------------------------
% "Синтез оптимального стохастического управления", 2011 год
% --------------------------------------------------------------------------------------------------



\chapter*{Заключение}
\addcontentsline{toc}{chapter}{Заключение}



В работе были исследованы вопросы синтеза оптимального регулятора с использованием принципа комбинированного управления, когда нейтрализовывалось влияние превалирующего возмущения (принцип компенсации), а побочные возмущения устранялись с использованием принципа обратной связи. Подробно исследовались вопросы синтеза оптимального регулятора с привлечением аппаратов динамического программирования; была установлена связь с принципами максимума; было показано, как свести методы к уравнению Беллмана. Полученные результаты удалось обобщить на случай многомерных стохастических систем и доказать принцип полной эквивалентности при синтезе оптимального регулятора для этого случая, был исследован вопрос стохастической устойчивости.

Результаты были применены к задаче синтеза управления для стабилизации полета баллистической ракеты в турбулентной атмосфере. Исследовались как вывод на <<номинальную>> траекторию, так и удержание курса на ней. Для данной задачи, используя теоретические результаты, были сконструированы алгоритмы нахождения оптимального управления, дано их математическое обоснование. Был также проведен анализ этих алгоритмов и установлены условия их применимости.

Кроме того, был разработан математический программный пакет, позволяющий программно смоделировать регулятор и использовать алгоритмы для проведения компьютерных экспериментов. Эти эксперименты не только подтвердили теоретические результаты, но и позволили обратить внимание на проблему, связанную с тем, что для наведения с использованием алгоритмов достаточно большую роль играет время наведения. Если время наведения окзывается недостаточным, то алгоритмы не только не позволяют решить поставленную задачу, но даже ухудшают результаты. Были даны численные оценки порогового значения времени наведения для конкретной задачи.


% список литературы
\addcontentsline{toc}{chapter}{Список использованной литературы}

\begin{thebibliography}{wlab}

% 2 раздел

\bibitem{YES}
\emph{Ширяев А.\,Н.}
\newblock Вероятность.
\newblock М.: Наука, 1980, 576 с.

\bibitem{BELLMAN1}
\emph{Беллман Р.}
\newblock Динамическое программирование.
\newblock М.: Изд-во иностранной лит., 1960, 472 с.

\bibitem{TCHAMRAN}
\emph{Tchamran A.}
\newblock On Bellman's Functional Equation and a Class of Time-Optimal Control Systems.
\newblock Journal of the Franklin Inst., 1965, vol.\,280, no.\,6, 493--505 pp.

\bibitem{FICHTENGOLZ}
\emph{Фихтенгольц Г.\,М.}
\newblock Основы математического анализа.
\newblock М.: Физматлит, 2008, т.\,2, 728 c.

\bibitem{OXENDAL}
\emph{Оксендаль Б.}
\newblock Стохастические дифференциальные уравнения
\newblock М.: АСТ, 2003, 408 с.

\bibitem{MERRIAM}
\emph{Мэрриэм Ч.\,В.}
\newblock Теория оптимизации и расчет систем управления с обратной связью
\newblock М.: Мир, 1967, 491 с.

\bibitem{KALMAN1}
\emph{Калман Р.\,Е.}
\newblock Когда линейная система управления является оптимальной?
\newblock Теоретические основы инженерных расчетов, 1964, т.\,86, №\,1, 69--84 с.

\newpage

\bibitem{XU}
\emph{Сю Д., Мейер А.}
\newblock Современная теория автоматического управления и её применение.
\newblock М.: Машиностроение, 1972, 544 с.

\bibitem{KRASOVSKY}
\emph{Красовский А.\,А.}
\newblock Системы автоматического управления полетом и их аналитическое конструирование.
\newblock М.: Наука, 1973, 560 с.

\bibitem{KALMAN2}
\emph{Kalman R.\,E.}
\newblock New methods and Results in Linear Prediction and Filterig Theory.
\newblock NY: John Wiley and Sons, 1963, 51 p.

\bibitem{BELLMAN2}
\emph{Беллман Р.}
\newblock Введение в теорию матриц.
\newblock М.: Наука, 1969, 702 с.

% 3 раздел

\bibitem{SKOROHOD}
\emph{Скороход А.\,В.}
\newblock Случайные процессы с независимыми приращениями.
\newblock М.: Наука, 1964, 280 с.

\bibitem{KLOEDEN}
\emph{Kloeden P.\,E., Platen E.}
\newblock Numerical solution of stochastic differential equations.
\newblock Berlin: Springer-Verlag, 1992, 632 p.

\bibitem{KUZNETCOV}
\emph{Кузнецов Д.\,Ф.}
\newblock Численное моделирование стохастических дифференциальных уравнений и стохастических интегралов.
\newblock СПб.: Наука, 1999, 459 с.

\bibitem{KUSHNER}
\emph{Кушнер Г.\,Дж.}
\newblock Стохастическая устойчивость и управление.
\newblock М.: Мир, 1969, 200 с.

\bibitem{HASMINSKI}
\emph{Hasminski R.\,Z.}
\newblock Stochastic Stability of Differential Equations.
\newblock Sijthoff & Noordhoff, Alphen naan den Rijn, 1980, 527 p.

\bibitem{ARNOLD}
\emph{Arnold L., Wihstutz V.}
\newblock Lyapunov exponents.
\newblock Springer Lecture Notes in Mathematics, 1986, Vol. 1186, 561--588 pp.

\bibitem{AUSLENDER}
\emph{Ауслендер Э.\,И., Мильштейн Г.\,Н.}
\newblock Асимптотические разложения показателя Ляпунова для линейных стохастических систем с малыми шумами.
\newblock Прикладная математика и механика, 1982, т.\,46, 3 вып., 358--365~c.

\end{thebibliography}




\end{document}
